\quellseite{331}
\chapter{§ 26. THE CATECHISM}

*G. v. Zezschwitz, ii 1, \emph{Der Katechismus}, 1864. K. Buchrucker, pp. 119-135. E. Sachsse, pp. 338-365. *F. Zange, pp. 73-142. A. Eckert, pp. 36-47; 53-55; 116-163. O. Baumgarten, pp. 68-72; 82-84. *J. Gottschick, pp. 142-158. J. Berndt, pp. 59-68; 99-105. R. Kabisch, pp. 137-139; 240-246. E. Chr. Achelis, pp. 420-443. A. Rude, pp. 100-117. Chr. Burkstuemmer, pp. 157-169. J. Steinbeck, pp. 201-218. K. Knoke, \emph{Ueber Katechismusunterricht}, 1886. G. Heimerdinger, \emph{Zur Reform des Katechismusunterrichts}, 1890. G. v. Rohden, \emph{Ueber christocentrische Behandlung des luth. Katechismus}, 1891. *L. Schultze, \emph{Katechetische Bausteine z. Religionsunterricht in Schule u. Kirche}, ⁵1891. A. Braasch, \emph{Reform d. Religionsunterrichts in d. Volksschule}, 1891. Malo, \emph{Zur Katechismusfrage gegen von Rohden}, 1892. W. Bornemann, \emph{Der zweite Artikel im Lutherschen kleinen Katechismus}, 1893. G. Voigt, \emph{Die Bedeutung des christl. Religionsunterrichts}, 1895. J. D. von der Heydt, \emph{Der Religionsunterricht in Schule u. Kirche}, 1896. *S. Bang, \emph{Katechetische Bausteine z. christocentrischen Behandlung d. I. Hauptstuecks}, ²1897. K. Just, \emph{Der abschliessende Katechismusunterricht}, 1896-07. A. Eckert, \emph{Die Formalstufen im Katechismusunterricht in Ev. Schulblatt}, 1898. F. H. Kähle, Der Kleine Katechismus L's., ¹⁹1902. G. v. Rohden, \emph{Ein Wort zur Katechismusfrage}, ³1902. *H. Brammer, \emph{Neue Bahnen f. d. Religionsunterricht}, 1900, 1903. R. Staude, \emph{Katechismusunterricht}, 1903. S. Bang, \emph{Zur Reform d. Katechismusunterrichts}, 1904. M. Reu, Explanation of Luther's Small Catechism, 1905. *A. Auffahrt, \emph{Die religiöse Frage und die Schule}, 1905. A. Reukauf, \emph{Didaktik d. ev. Religionsunterrichts in d. Volksschule}, 1906. *R. Staude, \emph{Das 8. Schuljahr}, ³1907. R. Hempel, \emph{Zum Katechismusunterricht}, ⁵1907. *H. Huebner, \emph{Wie laesst sich der Katechismusunterricht einfach, interessant und fruchtbar gestalten}, 1907. *H. Huebner, \emph{Was der kleine Catechismus fuer ein grosser Schatz ist}, 1907. A. Richter, \emph{Religionsunterricht oder nicht}, 1908. Arzt, \emph{Welche Maengel zeigt der gegenwaertige Religionsunterricht und auf welche Weise ist ihnen zu begegnen}, 1908. *W. Rein, \emph{Stimmen zur Reform des Religionsunterrichts}, 1904-1908. \emph{Die Umgestaltung des Religionsunterrichts i. d. saechsischen Volksschulen}, ed. by the Vorstand d. saechs. Lehrervereins, 1908. Pfeiffer, \emph{Der Sturmlauf gegen die Zwickauer Thesen}, 1909. Th. Franke, \emph{Der Kampf um den}
\quellseite{332}
\emph{Religionsunterricht}, 1909. *S. Bang, \emph{Luthers Kleiner Katechismus ein Kleinod der Volksschule}, 1909. *J. W. Horine, \emph{The Catechist’s Handbook}, 1909. *H. Matthes, \emph{Neue Bahnen f. d. Unterricht in Luth. Katechismus}, 1909. *F. Rendorff, \emph{Das Problem der Konfirmation und des Religionsunterrichts in d. Volksschule}, 1910. G. Rietschel, \emph{Zur Reform d. Religionsunterrichts in der Volksschule}, 1909. E. Sachsse, \emph{Zur Reform des Religionsunterrichts in der Volksschule}, 1909. F. Niebergall, \emph{Biblische Geschichte, Katechismus. Gesangbuch}, 1910. H. Matthes, \emph{Der Religionsunterricht im Dienste d. Erziehung innerhalb d. relig. Gemeinschaft}, 1910. *F. Eberhard, \emph{Der Katechismus als paedagogisches Problem}, 1912. *B. Doerries, \emph{Erklaerung d. kl. Katechismus. Ein Beitrag zur Reform d. Religionsunterrichts: Zehn Gebote}, 1912; \emph{Der Glaube}, 1912. *K. Eger, \emph{Evang. Jugendlehre}, 1912. *Th. Kaftan, \emph{Auslegung d. luth. Katechismus}, 1913. M. Reu, \emph{Die grosse Bedeutung des kleinen Katechismus Luthers}, 1913. *L. Himmels, \emph{Der Katechismus als Lebensbuch}, 1922. G. Haccius, \emph{Der Kleine Katechismus in der evang. Heidenmission}, 1923. M. Reu, \emph{A New English Translation of the Small Catechism}, 1926.

If Biblical History is properly taught it will not only introduce the child into the past, but also exhibit to him the norms of a true Christian life in the present, that is to say, it will touch upon most of the subjects expressed in the Catechism. Notwithstanding, it would be a grievous mistake to regard specific instruction in the Catechism as superfluous. On the contrary, such instruction is indispensable for the reason that the truths diffused through Biblical History and gained only by means of inference, require to be combined into groups and to be turned to account from one and the same point of view. By removing the Catechism as an independent factor, we would pour contempt upon the historical development as it has taken place under the guidance of the Holy Spirit; we would disparage the Church’s catechetical labors of almost two thousand years; we would undervalue the revelation through the Word and those parts of Scripture, particularly of the New Testament, which are not strictly historical; last but not least, we should commit a pedagogic blunder of farreaching consequences because the youth of the Church as well as the
\quellseite{333}
common people require brief, compact sentences in which the most important elements of what they may and should believe are summarized, and which may serve at any time as a medium of orientation. This much is certain: if there were no Catechism, one would have to be created.

But that does not dispose of the question whether Luther’s Small Catechism still satisfies all the requirements of today. In our time this is largely denied with great vehemence. It is gladly conceded that, in the past, it was entitled to great consideration, and that it has left many blessings in its wake. For this reason there is no objection to recognizing it as a significant historic monument of the Evangelical Church. But as medium of instruction for the young it is either decisively repudiated or largely emptied of its contents. The misgivings felt as to the employment of Luther’s Catechism in the instruction of the young are both dogmatical and pedagogic. So far as the dogmatical objections pertain to the Evangelical faith as expressed in the Catechism, we pass them by; for the opponents have gradually acquired sufficient honesty to admit that the expression of the Evangelical faith as found in the Catechism tallies perfectly with that of the New Testament. Now, since that is the case and since this faith is laid down in Luther’s Catechism in simple form, unburdened with theological problems and terms (cf. pp. 110-113), Luther’s Catechism is fully adequate; we desire, neither for our own person nor for the instruction of the adolescent youth, any other expression of Christianity than that laid down in the New Testament. But it remains to examine as to their correctness and bearings those objections against the Small Catechism which are raised from the pedagogical point of view and, likewise, those dogmatical objections which are not based on opposition against Holy Writ.

We concede from the outset that the Small Catechism has often been handled, and is still handled, in such a way as to strike all healthy pedagogic requirements in the face. When children are made to commit
\quellseite{334}
it to memory at a stage of life when they are not yet equal to the task; when the teacher assigns it for memorizing without in the very least paving the way for its appropriation; when its material is not properly outlined and the teacher fails to show his pupils the inner structure of the component parts; when the explanation fails to shed the needed light upon it from the characters of Sacred History and to connect it with the fresh, throbbing life of the present; when the catechist becomes a lecturer and possibly even lectures beyond the scope of the catechism; when he never lets the children feel the fervor of his own devotion; when by adducing features from the history of salvation he attempts to make it something that it was never designed to be; when he enlarges it into a textbook of dogmatics and chokes its own fresh life—in all such cases the hour for catechization may become one of positive torment for the children, and the bond between the Church and her young may, in consequence, be loosed rather than strengthened. But is Luther’s Catechism, this “\emph{golden gem}”, to be made responsible for these disastrous results when it is maltreated and abused? Is the Catechism to be blamed when it is accorded a treatment radically in conflict with its unique character?

Two objections are directed specifically against the \emph{First Chief Part}. It is said (a) that it is not suited for the evangelical instruction of youth because it incorporates the Mosaic Decalogue, that it leads the children back to a sub-Christian level, and that it teaches Old Testament, instead of New Testament, ethics. But the objection is not valid: the Catechism is made to shoulder a reproach deserved by its interpreters. Luther’s explanation has so thoroughly infused \emph{New Testament life} into the Old Testament commandments as to constitute this one of the greatest advantages of his Catechism. This matter has been treated at length on previous pages (101-103); only one fact is to be added here owing to the present form of the \emph{First Chief Part}. The introductory words, “I the Lord thy God,” were omitted in the editions of the Catechism published in Luther’s life time; but through the influence of the Nuremberg Sermons for Children they have been incorporated in the Catechism and retained. Thereby the Evangelical character of this part has become even more pronounced. This quality, indeed, is not brought out when one departs so far from a correct understanding of the text in question as to endeavor to educe theoretically from the words “I am” the existence and personality of God, and from the words “the Lord thy God”, his essence and character. For such a teacher, these words are merely an opportunity for an extremely unfruitful,
\quellseite{335}
because altogether unchildlike, and, at least in this connection, perfectly superfluous and fatuous treatment which, while it may possibly enable the catechist to subject his dogmatical knowledge to a wholesome review, is positively discouraging to the souls of the poor children. In truth, the words of introduction invite and attract like the open door of home. For when God gives His commandments the superscription, “\emph{I the Lord},” or “\emph{I, Yahweh, am thy God}”, He opens therewith His heart and home and offers Himself to us with all His saving power and mercy. He does not begin with the cold, severe, difficult demand: “You shall fear and love Me”; but He begins with the gift of Himself to us, with a revelation of Himself as \emph{Saviour and Father}. Not until He has given Himself to us as God and Father and opened wide the gate of home, does He open His lips to pronounce the rules and regulations of His house in the observance of which His children are to exemplify their love. With His love He intends to arouse our love so that He may reap where He first has sown. Is it not the foundation of all true godliness that God gives Himself to us? And, again, could there be a more perilous and fatal error than the opinion that one must make his own way to God by means of legal obedience? We have therefore reason to be grateful that through these words the gospel has been given precedence in the \emph{First Chief Part}; the introduction reminds us of Baptism: Yahweh has become our God and Father, in the commandments He gives His children the regulations of His house, and true children of His will gladly observe them. The authority divine is not at all weakened by this treatment, but rather made more effective. From what has been said, it is evident that the children, in the \emph{First Chief Part}, are not led into a world altogether foreign to them, as has been averred, a world for which they lack every connecting link in their own life. On the contrary, just as in their home life, the father, author of their life and provider of their bread, enacts certain regulations to which they are required to conform, so the Father in heaven with the words of the introduction opens the gate of His home, whereupon He lays down the several rules of His house. These are things quite suited for the concept sphere of the child. Nor is (b) the other objection to the \emph{First Chief Part} well grounded that the explanation of the several commandments dovetails inadequately with the life of the present. When enumerating the various virtues and faults, Luther merely means to give examples of the transgression or fulfillment of the commandments, as is clearly indicated by the “etc.” added in the original editions in the explanation of the Ninth Commandment.
\quellseite{336}
Moreover, in the Large Catechism, Luther expressly said that those features which were most commonly disregarded should receive the most careful attention. According to the Reformer’s own intention, therefore, there is no reason whatever not to supplement the examples of true obedience to the divine will given by Luther through others taken from the life of today. Personally, we are of the opinion that virtually all that is necessary is already found in the original outline. What is needed is not so much new additions as an exhaustive amplification of what is given.

When it is asserted that the \emph{Second Chief Part} (a) postulates so profound an insight into sin and a faith so strong as altogether to transcend the experience of the child, the fact is overlooked that according to Luther’s intention, it is primarily the father or the catechist who confesses himself to be a lost and condemned sinner and expresses his faith. This is irrefutably proved in the explanation of the \emph{First Article} where “wife and child, house and home” are enumerated. This is further corroborated by the practice, in Luther’s day, of the father or the catechist pronouncing, and the children and servants repeating, the words. Moreover, it must be borne in mind that the initial stages of contrition and faith can well be wrought by the Holy Spirit in the period of childhood and that, during adolescence, they sometimes reach a surprising depth (pp. 262, 268). (b) When the assertion is made that the notion of creation ex nihilo, as contained in the \emph{First Article}, is not an essential part of a saving evangelical faith, a grave error is committed. For it is the indispensable foundation for the article that God is the absolute Lord of the world and that nothing can harm us save by His permission. That this conviction is essential to the Christian faith, no one can deny. It is sometimes said (c) that the explanation of the \emph{First Article} does not correspond to the conditions in which most of our children are reared because we find enumerated “house and home, wife and children, cattle and all goods”. Now, it is quite true that Luther when he wrote this explanation had in mind the rural family; but the fact must not be overlooked that this specification is merely intended as an illustration. It is not likely that serious difficulties will arise for an experienced catechist in this connection. If (d) the fact is emphasized that the socalled intermediate causes have been completely ignored, everything being traced directly back to God, it should be remembered that the Bible deals with the subject in precisely the same way without thereby denying the operation of the intermediate causes. The judgments expressed at this place are
\quellseite{337}
judgments of faith; and faith does not deal with powers and objects belonging to the economy of nature. It is rather the very nature of faith to look beyond the material world for the invisible one lying behind it; nor is it satisfied until in everything pertaining to our natural life, the concealed hand of God has been seized which through those selfsame intermediate causes supplies food and everything else that is needed for our body and life, protects us against all danger, and guards and keeps us from all evil. It is just such emphasis upon God as the \emph{causa prima} which is bitterly needed in our day where the highly cultivated natural sciences tend to blind men against the governing hands of the invisible God. Again, (e) if one objects to the fullness of assurance with which Luther confesses that “He richly and daily provides me with all that I need for this body and life, protects me against all danger, and guards and keeps me from all evil”, we reply (1) that it is a judgment of faith which is here expressed, viz., the confidence that one is safe in the arms of God for all time to come; (2) that the accent lies on “God” in contrast to one’s own strength and wisdom; and (3) that the terms “danger” and “evil” are not to be interpreted from the narrow standpoint of our limited human reason, but, in the light of Rom. 8:28, from that of the all-seeing divine ruler of the world and of every individual life. Offense has also been taken (f) because the Apostolic Symbol silently passes over all intermediate steps between the birth of Christ and His passion and also over the prophetic ministry of Christ, and because the same omission occurs in Luther’s explanation. But let it be remembered that this is precisely the method of Paul in all his letters, and that the passion and death of Christ is in reality the climax and crowning consummation of His work.

The objection against the \emph{Third Chief Part} is that the tenor of the explanation is too heavy. We readily concede that due to the lengthy sentences and massive thoughts there is an occasional lack of lucidity and that the catechist faces peculiar difficulties on this account. If, however, he has faithfully explained the first two Chief Parts, the children are able to comprehend what Luther offers in his explanation of the Third. The last trace of difficulty will disappear if he will direct the pupils to repeat in brief sentences the respective burden of the individual petitions (cf. our Explanation of the Catechism). It is also an easy matter to add the thought of missions in the Second Petition.

When it is said that the Fourth Chief Part is too largely oriented
\quellseite{338}
upon adult baptism and that infant baptism is virtually passed by, we answer that whatever pertains to the essence of Baptism is of necessity true of infant baptism as well if the latter really is Baptism.

While it is true, finally, that Luther erred in connection with the \emph{Fifth Chief Part} in that he applied the words of institution, “Given and shed for you for the remission of sins”, to the act of distribution, and not, as is proper, to the act of Christ’s death, and then, acting on such erroneous interpretation, fixed these words as central\footnote{Toward the end of his life Luther conceded their bearing upon the death of Christ without, however, casting aside their bearing upon the Sacrament; cf. Mathesius’ xv. sermon on the life of Luther.}, we do not endorse his error; but the catechist who is convinced that the sacramental blessing, the remission of sins, is bestowed not through the body and blood of Christ, but through the Word used in connection with the Sacrament, need not consider this formula—thrice recurring—as anything but recapitulation of the words of institution whereupon he will find his difficulty vanished. While the matter is not so simple for one to whom the body and blood of Christ is intended not only to seal but also to transmit the sacramental blessing of the remission of sins, the difficulty can be surmounted nevertheless. It would be altogether wrong for such reasons to take offense at Luther’s Catechism and to doubt its availability as a medium of instruction in the present. The few occasional difficulties are more than outweighed by the numerous merits of this peerless little book.

Accordingly, if Luther’s Small Catechism is the best summary of that faith which is the life of the mature congregation, pains must be taken in catechetical instruction not to obliterate its great merits, but rather to give them full scope so that the uniqueness of the Catechism is reflected in the instruction based upon it as its text. This will be accomplished most readily when instead of going elsewhere for leading points and grafting them as well or ill as possible upon the several parts of the Catechism, all of the instruction is nothing more or less than an \emph{unfolding of the fullness of the evangelical life of faith expressed in the words of the Catechism}. It was Loehe particularly who reaffirmed this principle; Theo. Kaftan in his \emph{Auslegung des Lutherschen Katechismus} (⁶1913), and likewise
\quellseite{339}
the \emph{Synodical Catechism of the Iowa Synod} have consistently carried it out.

This principle recognized, it follows naturally that instruction in the Catechism should be restricted to the truths that constitute the center of the Christian life. That involves the exclusion of special complementary parts which, it is thought, form a necessary addition to Luther’s Catechism, as for instance, a special section relating to the Word of God or the new life of the Christian (Beatitudes, Twin command of love). It involves also the exclusion of supplementary sections taken from dogmatics and sacred history whereby the verbal limits of the Catechism are transcended. There are catechists who believe a discussion of such concepts as revelation, religion, conscience, to be necessary to an introduction to the Catechism. They believe that no adequate treatment of the Second Chief Part is possible without first having laid down a general doctrine of God, His essence and attributes, and the Holy Trinity. They cannot refrain from interweaving the whole history and doctrine of creation although now separate courses in Biblical History are found everywhere; from making an “excursus” to the subject of good and bad angels; from foisting upon their treatment of the First Article the dogmatical division “creation, preservation, and government”. They consider as a weakness of Luther’s explanation of the creed his failure to interlink between the First and the Second Article the doctrine of the fall of man, and try to improve this condition by adding an elaborate statement concerning sin and the need of redemption. In connection with the Second Article they demand a supplement treating of the sinlessness of Christ or of His whole ministry. They think that the explanation of the Second Article must be treated according to the conventional scheme of dogmatics, “Christ’s Person and Work”. It is simply incomprehensible to them how one can dispense with so important(!) a distinction as that of the prophetic, priestly,
\quellseite{340}
and royal offices of Christ (a distinction which, as a matter of fact, is to be made only with due caution, and which was not fully developed until the seventeenth century). And then, to cap the climax, they introduce an extended discussion on the Two States of Christ (although they have treated the same material before, even if in different connection!). In spite of the fact that the expression “redemption from death” is found midway between redemption from sin and redemption from the power of the devil, they cannot refrain from adding the distinction—quite appropriate, to be sure, at the right place—in regard to natural, spiritual, and eternal death. In connection with the Third Article, they endeavor to give a detailed explanation of the person of the Holy Spirit; and they believe themselves guilty of a serious omission when they fail to enumerate all the steps of the order of salvation ever constructed in dogmatics or to force everything that can be said on the subject into the few words found in the catechetical text proper. Instead of demonstrating in the light of the Lord’s Prayer itself what true prayer is, they premise an elaborate doctrine of prayer. The conception of the Sacrament with which they operate has a cast altogether dogmatical; and then they strain under obstacles self-imposed when they attempt correctly to state the heavenly blessing of Baptism, etc. That these elucidations are correct in themselves, is not the point in question. They are not in place in any kind of catechetical instruction which would preserve the uniqueness of Luther’s Catechism. That book excludes the technical ballast of dogmatical terms as much as is practicable; and instead of serving as an introduction to theology, it has no other purpose than to serve as a succinct and lucid expression of the substance of the evangelical life of faith.

If instruction in the Catechism is confined to what is central in the Christian life, it will, naturally, refrain from fusing the Five Chief Parts into a system. Such attempts
\quellseite{341}
have been made again and again. The several parts have been joined into a connected whole by such an arrangement as this: The First Part is to show us the disease, the Second the physician, the Third the means whereby the physician is to be summoned, the Fourth and the Fifth the remedies employed by the physician to effect a cure. Or by the following arrangement: The First Chief Part shows us the goal to be reached, the Second shows us faith as the way to this goal, the Third, Fourth, and Fifth show how the power to prosecute the way is to be found, namely in prayer and the Sacraments. According to the first systematic effort here described, the Ten Commandments are of moment exclusively as a mirror for sin, according to the second as the rule and norm of the new life. The attempt is made to justify both these conceptions by appealing to Luther, but without warrant (cf. pp. 111 f.); for, according to the Small Catechism at least, the First Chief Part is, indeed, primarily the norm of the true Christian life, but also the mirror for sin (see the part treating of Confession). In a type of instruction in the Catechism that takes its cue from the Small Catechism itself, the Decalogue will be used to impress the ideal Christian conduct upon the vision of the pupil with as much vividness and animation as possible. Through such a treatment, the Decalogue, apprehended in an evangelical sense, will of itself become a mirror for sin. Thus the danger of formulating a system has been obviated; for that can exist only when the Decalogue is treated either as a mirror for sin or as the norm of the new life, but not when it is treated as both. At the same time a service is rendered to the young who often have been tormented with questions relating to the connection between the one Chief Part and the other. The young and the common people require brief and definite statements as a treasure secure and rich, but not a system. — Moreover, the treatment of the
\quellseite{342}
Catechism is rather one-sided when it is based upon the notion that each Chief Part contains but one side of Christianity and excludes the others. The fact is that each Chief Part contains all of Christianity although the light and the relation in which it is viewed do not remain the same. In the First Chief Part Christianity is viewed as a divine requirement; in the Second, as the expression of personal conviction; in the Third as a divine gift to be received through prayer; in the Sacraments as salvation appropriated by visible acts; in the additions, as exemplified in the various relations of life.

Where the principle stated above is adhered to, a third wrong method stands excluded. We refer to the method which is not satisfied with merely unfolding the explanation of Luther but insists upon adding another explanation of the Catechism text proper (here and there, even one that is not at all in concord with that given by Luther). Such method has a most disturbing effect when, e.g., each step of the humiliation and exaltation of Christ is examined separately, and its meaning and importance, as distinguished from that of the preceding and following steps, is brought out; or when in connection with the Third Article its five component parts are treated separately (e.g., in \emph{Caspari’s explanation}). Such method, considered from the standpoint of pedagogics, is thoroughly vicious, connoting as it does a misapprehension of one of the greatest merits of the Catechism (cf. pp. 103 ff.). No, the text of the Catechism has been explained by Luther in accordance with Scripture, and therefore in an authoritative manner. For this reason the text confronts us with no other requirement than that a brief definition be given of difficult words; that when occasion suggests it, a few historical remarks be interspersed; and that wherever necessary, the outline of the text be disclosed, as in the Second and Third Articles, the catechist always
\quellseite{343}
taking his cue from Luther’s appended explanation (cf. Reu, \emph{Explanation of Luther’s Small Catechism}, pp. 63 f., 73-76, 88 f.). Instead of introducing the pupil to theology and dogmatics, the catechist should make every effort to bring out the very meaning of the truly central religious elements, to impress it upon the hearts of the children, and to make the application to their life. He will try to eradicate the erroneous notion as though the ‘lifeless, useless dogmas’ of the Catechism must simply be hoarded up in the memory; wherever possible he will rather exhibit the great significance, for the present life of the pupil, of the several truths of the Catechism. It will be necessary for him to draw upon Biblical History, upon secular history, upon the history of the Church and of missions; if need be, also upon the kingdom of nature, in order to gain an appropriate starting-point or illustrative material. He will find much useful material in H. Caspari, \emph{Geistliches und Weltliches}, ²1915; J. H. A. Fricke, \emph{Handbuch des Katechismusunterrichts, zugleich Buch der Beispiele}, ²1892; L. Pestalozzi, \emph{Die christliche Lehre in Beispielen}, ³1901; P. von Zychlinski, \emph{Illustrirende Aussprueche, Sentenzen und Geschichten zu Gottes Wort}, 1900; J. Besch, \emph{Aus der Lernstube des Lebens}, ²1913; and \emph{Narratives on the Catechism}, Columbus, O.; valuable sketches and narratives from the field of missions are found in Warneck, \emph{Die Mission in der Schule}, ¹⁴1912; Schaefer, \emph{Die Innere Mission in der Schule}, ⁷1912; and Koenig, \emph{Die Mission im Katechismusunterricht}, 1913. But whatever material the catechist may have selected for the purpose of bringing home to the pupil the truth in a quickening and touching manner, of impressing upon him its value, and of producing permanent impressions and promptings of the will—always the statements of the Catechism must constitute the cord upon which everything is strung, or the ultimate issue of the catechesis; nor should any element of knowledge receive consideration
\quellseite{344}
which is not an organic outgrowth of the Catechism. Only in this way will the youth of the Church be well grounded in the Catechism, katechismusfest; only in this way will it become an instrument for the orientation of their life in days to come. The teacher who is not content with letting his explanation turn upon Luther’s explanation, might well be asked why he insists at all upon an exact memorization and repeated reviews of the same.

Finally, all instruction in the Catechism should be characterized by the tendency — manifested also in Luther’s Enchiridion — to individualize instead of dealing in impersonal generalities (compare “we should” in the First Chief Part; “I” and “me” in the Second Chief Part; “we,” “us,” “our” in the Third Chief Part). The teacher should be a \emph{Seelsorger} rather than a dogmatician: he should try to arouse not only the intellectual faculties, but also the feelings and the will, in short the whole man. To summarize: instruction in the Catechism should possess all those marks that constitute the superlative excellence of the Small Catechism. The more it preserves and reflects its unique character, the more likely it will become truly evangelical and pedagogically correct. And it will reflect its uniqueness best when the instruction from the beginning has no other aim than to be an unfolding of the fullness of evangelical faith, as it is expressed in the words of Luther’s Catechism.

We take the liberty of adding a few hints as to a correct understanding of the Small Catechism. While we must confine ourselves to essentials at this place, we refer for details to M. Reu, \emph{Explanation of Luther’s Small Catechism} (larger edition, Chicago, ¹¹1926; smaller edition, ibid., ³1926); M. Reu, \emph{A New English Translation of the Small Catechism}, 1926, to Vol. xxx of the \emph{Weimar Edition of Luther’s Works} Weimar. Boehlau, 1910, upon which O. Albrecht has bestowed great care; A. Ebeling, \emph{Historisch-kritische Ausgabe von Luthers Kl. Katechismus}, Hannover, 1901; A. Hardeland, \emph{Luthers Katechismusgedanken in ihrer Entwickelung bis zum Jahre 1529}, Guetersloh, 1913; J. Meyer, \emph{Luthers Grosser Katechismus, mit Kennzeichnung seiner Predigtgrundlagen},
\quellseite{345}
Leipzig, 1914; and esp. to Th. Kaftan, \emph{Auslegung des Lutherischen Katechismus}, Schleswig, 1913, the best commentary to Luther’s “golden gem”. In the last named work all important books on the subject are noted. Parts of Kaftan’s book are extant also in the English language, in the form of \emph{The Catechist’s Handbook}, by J. W. Horine, Philadelphia, 1909. While my own views in the premises are largely in agreement with those of Kaftan, justice to our subject requires the statement that, as my notes show, my views, on the whole, were formed years before I became acquainted with the book of Kaftan. The impulse for my conception of the Catechism in general came from the foreword of W. Loehe’s Catechism of 1845.—Abundant material relative to the history of the exposition of the Catechism is found in M. Reu, \emph{Quellen zur Geschichte des kirchlichen Unterrichts im evangelischen Deutschland zwischen 1530 und 1600}, Guetersloh, 1904 ff.

\emph{The introduction should be brief.} Following Brenz’s example (a similar procedure is found already in \emph{Fundamentum aeternum felicitatis seu summa religionis Christianae} which appeared in print in 1498 after some years of manuscript circulation) and that of many expositors of Luther’s Catechism, notably Bischoff (1599; cf. Reu i, 2, pp. 222 f.), the teacher will use as a starting point the sacred act by which the catechumens have become children of God, i.e., Baptism; touch upon Holy Scripture which bears witness to that fact; and then pass on to the Catechism which, in five phases, is designed to teach them in what the faith and the life of God’s true children consist.

\section{The First Chief Part}

“The Ten Commandments,” Luther says at one place, “reveal to us what we have been, what we are, and what we should be.” His mode of presenting the subject indicates that he wishes the last point to be stressed: always he begins his explanation with the words “we should”. His aim is to hold before the young the ideal Christian life and to show them, in clear outline, what they should do and what they should avoid. But are the commandments of Sinai adapted to that purpose? Are they not an integral part of the Law which has found its end in Jesus Christ, in fact, are they not that very Law summarized? Quite true, in their Old Testament form they are obsolete as is clearly seen when we compare the Fourth and the Third Commandments with Col. 2: 16 f. and Gal. 4: 10 f. Nor did Luther incorporate them in the Catechism in their Old Testament form; still less did he give them an Old Testament interpretation (cf. pp. 101-103). But he gathered the divine thoughts of eternal value that were contained in the Old Testament form and these he explained in the sense
\quellseite{346}
and spirit of the New Testament. This, and this alone, is the acceptation of the Ten Commandments in the present: so understood, they are gratefully welcomed by the Christian, eager to do the will of God of which he is often ignorant, and gladly accepted as the rule and norm of his life. The children will undoubtedly observe the difference between the biblical wording of the Decalogue (as found also in their Biblical History) and that in the Catechism; this gives the catechist a fine opportunity to spread healthy evangelical conceptions of Holy Scripture. In fact, it is his duty to make this matter plain at the outset because the Reformed never cease accusing Luther of mutilating the text of the Bible in this respect. Afterward, in connection with the Third Commandment, a similar opportunity presents itself.

The Commandments are divided into three parts: (1) The Introduction; (2) The Commandments proper; and (3) The Conclusion. \emph{The Introduction}—omitted by Luther who followed the traditional text—was rightly restored almost everywhere (cf. p. 312 f.). Luther himself took note of it not only in his sermons on Exodus (1525) and Deuteronomy (1529), but also in his catechetical sermons of 1528. The exodus from Egypt referred to in Ex. 20 and Deut. 5 was to Luther the “sign by which the Jews were to lay hold of God”; what the exodus was for the Jew, redemption by Jesus Christ is for the Christian; at one place there is even a reference to Baptism (Weimar Ed. xvi, p. 425). “For this reason it is not necessary that they [the Christians] should fear and love God under the name of deliverer from Egypt: He has received another name now, namely Christ who, true God, has died for us” (Weimar Ed. xxviii, p. 605). In his annotations on the Decalogue of 1530, he therefore is able to say in regard to the words of introduction: \emph{Promissio omnium promissionum, fons et omnis religionis et sapientiae caput, evangelium Christum promissum complectens} (Weimar Ed. xxx, 2, p. 358). Accordingly, the catechist is following in Luther’s steps when he uses the Introduction as an occasion to remind his pupils of their redemption by Jesus Christ or of the baptismal sacrament by which the former is appropriated, and says to his class of Christian children: “These commandments are given you by Him, who, through Baptism, has become your God and Father. But if He has granted you this wonderful blessing that He has become your God and Father, it is His right to give you rules and regulations for your conduct, and He can justly expect you to live up to them in grateful love.” Sometimes a passage from Luther’s Table Talk (Erlangen Ed. lviii, p. 206 f.) is adduced as disproving the statements made above;
\quellseite{347}
but the objection is not well taken: Luther does not argue at all against the retention of the Introduction, but against its inclusion among the commandments (cf. Reu i 1, p. 451, 758), because to him such a procedure would, naturally, be a perversion of a “promise” into a commandment, a suppression of the comforting gospel which stands at the head of the commandments (cf. p. 313).

We are ignorant of the arrangement of the Commandments upon the tables of Moses. This should prompt the catechist to pay no attention to this point when he discusses the question as to the division of the commandments; still less will he exaggerate the division 1-3 and 4-10 into a shibboleth of the Lutheran Church. While there is no adequate reason to depart from this division, no attempt should be made to maintain, with an unwarranted reference to Matt. 22:36-40, that the first three commandments treat of that which we owe to God while the last seven treat of that which we owe to our neighbor; for the child will not number his parents among the “neighbors” (note that Luther does not introduce the term “neighbor” until he comes to the Fifth Commandment). For this reason Melanchthon and many interpreters of Luther’s Catechism made a distinction between what we owe to God and what we owe to men. While this distinction has a better warrant, it is still misleading. A mere glance at Luther’s explanation is sufficient to assure us that our duty toward God is the subject under consideration also in the last seven commandments; the difference is this that in the first group the obligation imposed is direct, while in the second group it is indirect. The first table shows what we owe to God himself, and the second table what we owe to our fellow-man for the sake of God. A reference to Matt. 25:40 will at once make the subject clear to the child: he will begin to comprehend the important truth that all he does bears not only on men, but primarily on God (Psalm 51:4).

He who has come to a proper understanding of the Introduction will find the First Commandment the necessary logical sequel. If God through Baptism has become our God, it is His right to require that we should have no other gods before Him, but have Him alone as our God. The question when that is the case that we have God as our God, is usually answered by Luther, quite in accordance with the New Testament, in that he simply points to faith. E.g., he writes in the explanation of the Decalogue as found in his Sermon on Good Works (1520): “Since I alone am God, thou shalt place all thy confidence, trust and faith in Me alone, and on no one else. For that is not to have a god, if you call him God only with your lips, or worship him with
\quellseite{348}
the knees or bodily gestures; but if you trust Him with the heart, and look to Him for all good, grace and favor, whether in works or suffering, in life or death, in joy or sorrow . . . And this faith, faithfulness, confidence deep in the heart, is the true fulfilling of the First Commandment; without this there is no other work that is able to satisfy this Commandment. And as this Commandment is the very first, highest, and best, from which all the others proceed, in which they exist, and by which they are directed and measured, so also its work, that is, the faith or confidence in God’s favor at all times, is the very first, highest, and best, from which all others must proceed, exist, remain, be directed, and measured. Compared with this, other works are just as if the other Commandments were without the First, and there were no God” (Weim. Ed. vi, p. 209). Similarly he says in his sermons on Exodus (1525): “All the First Commandment requires is genuine faith and trust in God; it does not enjoin anything external. But no one can fulfill it unless the Holy Spirit work it in his heart” (Weim. Ed. xvi, p. 445; cf. p. 464 where he repeats the exposition of 1520 almost literally). In the same strain he expresses himself once more in his Catechetical Sermons of 1528: “\emph{Ergo intentio huius praecepti est dass es will gebieten einen rechten Glauben}” (W. E. xxx 1, p. 28), and in the Large Catechism: “Therefore it is the intent of this commandment to require true faith and trust of the heart which settles upon the only true God, and clings to Him alone (W. E. xxx 1, p. 133; Concordia Triglotta, p. 581, § 4). Just because it requires faith, he aptly calls the First Commandment “the chief part of all our Christianity” (in his sermons on Deuteronomy of 1529; W. E. xxviii, p. 601), and properly describes the Church as “the number of those who trust in nothing save the mercy of God alone and understand the First Commandment” (W. E. xxviii, p. 580). In the sketch De Loco Iustificationis he even writes: “pueri et infantes confirmant suo catechismo solam fidem absque operibus iustificare . . . Primum praeceptum est promissio quod velit esse Deus. Et fidem exigit ante omnia opera quae sequentibus praeceptis exiguntur. At fide secundum primum praeceptum habita filii Dei sumus, remissis iam peccatis, ipsa fide iusti” (W. E. xx p. 663). Now, when Luther in his Small Catechism in answering the question when we regard God as our God, points to something threefold, it is not likely that we are to look for anything else in his answer than a development of the conception of faith from the pedagogical standpoint. And indeed, he who believes in God as his God and Father is bound to fear, love, and trust in Him above all things. To trust in
\quellseite{349}
him; for, to trust in God as One who in Jesus Christ has become our Father and who now so governs all things that nothing can harm us against His will,—that is the very central part of our faith. To love Him; for, to love means to surrender to God; and this, too, is involved in faith: at first (logically) faith can, indeed, do no more than receive God and His precious gift, but on the strength of these gifts divine it subsequently rises to personal self-surrender. To fear Him; for how should faith maintain itself without feeling awe of that God who has shown Himself so great and sublime in all His deeds? The assertion has indeed been made, in view of the conclusion of the Ten Commandments, that Luther did not have in mind the awe of God natural to a believer, but the fear of God's wrath and punishment. Now it is quite true that Luther when he has occasion to explain the expression “jealous God” does stress this aspect; and that the fear of God's wrath is for the Christian as long as he is garbed in the flesh, and just because of that fact, an additional motive for obedience, has been conceded also by Luther. We read, e.g., in a letter dated Oct. 27, 1527: “\emph{Timeant poenam et infernum omnes impii, Deus aderit suis, ut simul timeant Deum cum poena. Neque fieri potest, ut sine timore poenae sit timor Dei in hac vita, sicut nec spiritus sine carne, etiamsi timor poenae sit inutilis sine timore Dei}”. (Enders vi, 109). But this does not at all justify the view that Luther, in his explanation of the First Commandment and in repeating the words, “We should fear and love God”, in his explanation of the other commandments, had this idea of fear in mind. This is impossible because his explanation of the First Commandment is always under the influence of the superscription which is to him sheer gospel and promise, and because he here unfolds the conception of faith. In the third series of Catechetical Sermons, of 1528, in which the First Commandment, even more emphatically than usually, is summarized as timor and fides, timere, according to one manuscript, is also found in juxtaposition with colere and venerari (Weim. Ed. xxx 1, p. 59, lines 32 and 34). In other words, what Luther desires to bring out forcefully is this: God is so great and august, so sublime and wonderful that, filled with holy awe, the Christian can only marvel and adore in His presence and that he must tremble at the very thought of ignoring His will or offending His love by ingratitude or sinful conduct. No one can doubt that this idea of God's sublimity, majesty, omniscience, omnipresence is a most powerful motive for a righteous life, and at the same time a motive most appropriate for him who, in faith, has received a new vision of
\quellseite{350}
God’s greatness and wondrousness. It is the conclusion of the commandments that supplies the thought of divine wrath and punishment; it reminds the Christian: If you, in your carnal mind, are so thoughtless and foolish that the idea of God’s wondrous greatness does not move you to sanctify your life, remember that this great and august God can and will prove His greatness also by punishing, in terrible wrath, all those who despise His commandments and frivolously or maliciously ignore His will!” According to Luther, then, the fear of God’s wrath and punishment is added—but only added—to the more important factor (awe of God’s greatness as inherent in faith and inseparable from it). Ten times he speaks of the latter, only once of the former; just in proportion as we can dispense with the former in our life and the latter becomes the dominant factor in it, we shall become more and more what we should be. The catechist who has recognized this fact will, therefore, with a reference to a passage like Gen. 17:1 or Tob. 4:6, explain this conception to his children in the following manner: “To fear God, means always to have before our eyes the great and majestic God, the almighty and omniscient One, so that we fear to sin against Him.” He had better not introduce the dogmatic distinction between filial and servile fear since the fear taught in the First Commandment, thus defined, does not quite tally with filial fear.—Likewise, since he expounds the First Commandment to Christian children, he will refrain from any reference to pagan idolatry; still less will he accept the more than questionable distinction between ‘coarse’ and ‘refined’ idolatry, but, with Luther (cf. the Large Catechism), confine himself to that idolatry which fears, loves, and trusts something else more than it does God, and which, for a Christian, to whom God has revealed Himself as his God and Father, is coarser and more wicked than any species of pagan idolatry. Considered pedagogically, a discussion of heathen idolatry at this point is a deflection involving the danger that the children regard Christian idolatry as the lesser offense. The vigorous underscoring of the “above all things” will so completely engage the teacher’s attention that no time is left for anything else. The best exposition of the First Commandment of the Lutheran Catechism is found in the Nuremberg Sermons for Children (in Reu i 1, p. 465 ff.), except that at one place a reference to God’s wrath has erroneously crept into the explanation of the fear of God.

In explaining the remaining commandments, the catechist will take pains regularly to stress the “we should fear and love God” of Luther lest he inculcate Old Testament instead of New Testament,
\quellseite{351}
medieval instead of evangelical and reformatory, ethics. Eventually he will accept no other answer to the question, What does God require in this or that commandment? than this: We should fear and love God. In this way he will effectively drive home both the unity of Christian ethics and the fundamental truth that no work pleases God unless it comes from a heart that fears and loves Him (p. 102 f.). Says Luther in the Large Catechism: “Thus the First Commandment is to shine and impart its splendor to all the others. Therefore must this declaration run through all the commandments, like a hoop in a wreath, joining the end to the beginning and holding them all together; that it be continually repeated and not forgotten, etc.” (Weim. Ed. xxx 1, p. 181; ‘Triglotta, p. 675, § 326-329). In his Sermon on Good Works (1520) Luther calls the faith required in the First Commandment the Christian’s “\emph{health}” and says: “. . . a man must first have \emph{health}, which will work all the works of all the members of the body. So faith also must be in all works the master-workman and captain, or they are nothing at all” (Weim. Ed. vi, p. 213). The other commandments, from the second to the tenth, merely bring out in what manner the fear and love of God is to show itself in thought, word, and deed: ‘whereunto it is to “\emph{urge and impel us}”’ (Weim. Ed. xxx 1, p. 181, line 6; cf. Reu, Explanation of Luther’s Catechism, p. 33). That Luther, in constantly repeating the First Commandment in his explanation of the others, does not likewise repeat the word “trust” but confines himself to the words “fear and love”, is probably explained by the fact that both the positive and negative elements in the several commandments correspond better to the words “fear and love” than to the word “trust”; especially since Luther used it in a unique sense—trusting God for leading us in the best way. This view is borne out by the Large Catechism; for there we read: “Whatever you lack of good things, expect it of Me, and look to Me for it, and whenever you suffer misfortune and distress, creep and cling to Me. I, yes I, will give you enough and help you out of every need; only let not your heart cleave to or rest in any other . . . . Look to Me for everything, and regard Me as the one who will help you and pour out upon you richly all good things . . . We are to trust in God alone, and look to Him and expect from Him naught but good, as from one who gives us body, life, food, drink, nourishment, \emph{health}, protection, peace, and all necessaries of both temporal and eternal things. He also preserves us from misfortune, and if any evil befall us, delivers and rescues us, so that it is God alone (as has been sufficiently said) from whom we receive all good, and by whom we are delivered
\quellseite{352}
from all evil” (Weim. Ed. xxx 1, p. 133 f.; Tr. p. 581 ff.). So understood, trust can not at once be recognized by the child as a motive for the fulfilment of the Law. The child can grasp this idea much better in the context in which Luther has it in the conclusion: If now we keep the commandments, and still perceive but little of the promised ‘grace and every blessing’, we should not be dismayed at this unexpected experience, but cling to our conviction and trust to God that, in due time, he shall acknowledge our actions as done for Him and crown our efforts with a happy issue.—That Luther meant to correlate the fear of God rather with the negative features of the commandments, and the love of God with the positive ones, is a view doubtless warranted by the facts; but that, after all, not much weight is to be attached to this distinction, is evident from the Sixth Commandment which, like the others begins with “We should fear and love God” and yet contains only positive statements.

The conjunction “dass” in the explanation of Commandments Two to Ten has not received the same acceptation everywhere. The first Latin translation (in the Enchiridion piorum precationum, 1529) and, in part, also the third (by Justus Jonas, 1539), change the whole construction and offer instead, e.g.: \emph{Debemus Deum timere et amare, divinum eius verbum non contemnere aut negligere, sed illud potius magnificare et venerari, libenter tum audire ab aliis, tum alios docere}. The second Latin edition (by Sauromannus, 1529) and the fourth (by Hiob Magdeburg, 1560) construe the conjunction as final, rendering it by ut, resp. ne (e.g. \emph{Debemus Deum timere et diligere ne divinos sermones, eius verbum contemnamus, sed ut sanctum repetemus} etc.; or, \emph{Timendus et amandus Deus est ne sacras conciones et ipsius verbum negligamus, sed ut illud sancte veneremur} etc.). The Greek translations have δεῖ with the infinitive (Hiob Magdeburg: Δεῖ ἐν φόβῳ καὶ ἀγάπῃ τοῦ θεοῦ ἡμᾶς μὴ καταφρονεῖν κηρύγματος τοῦ θείου λόγου κτλ.) or they have ὥστε or ἵνα. Menius (1532) and others think that it introduces the manner in which we are to fear and love God (to make the matter plainer they add also; e.g., \emph{Wir sollen Gott also fuerchten und lieben, dass wir die Predigt und sein Wort} etc.). With this corresponds the former English translation: We should so fear and love God as not to despise His word, etc. It appears that Luther looked upon the conjunction in question as introducing the contemplated result, for he writes in the Large Catechism: “\emph{In like manner shall such fear, love, and trust urge and impel us not to despise his Word, but to gladly hear, learn and honor it, and esteem it holy}.” This quotation warrants the conclusion
\quellseite{353}
drawn by Albrecht that the clauses beginning with \emph{sondern} are to be construed as governed by \emph{dass}, and not as independent clauses. The explanation of the Sixth Commandment does not contradict this view; and the explanation of the Eighth Commandment which repeats \emph{sollen} after \emph{sondern} represents an exception.

The history of Luther’s Catechism in the English language evidences the difficulty to reconcile the meaning couched in Luther’s original with the genius of the English language. The following quotations may serve as illustration: We should fear and love God that we may not despise preaching and His Word; We should so fear and love God as not to despise preaching and His Word; We should fear and love God and for His sake not despise. . .; We should fear and love God and not despise preaching and His Word. From the standpoint of diction the last named translation deserves preference. Although syntactically the two clauses of the sentence are coordinate, a brief and simple explanation is sufficient to bring out the original meaning of Luther’s.

As regards the text of the Second Commandment, the threat connected with it in the Old Testament was not accepted by Luther nor is it found in the editions of the Catechism which appeared during his life-time. The catechist has so much less reason for wishing to correct Luther as the demand for the addition in this country is largely based upon the erroneous conception, that we are tied to the words of the Old Testament (cf. p. 323 f.). At all events, the catechist will take care not to interpret Luther’s “fear” in the light of the divinely threatened penalty. By that construction he would merely vitiate what he has taught concerning the meaning of this word in the First Commandment and what he will teach again in Commandments Three to Ten. Moreover, he would anticipate the comments to be made, according to Luther, upon the Conclusion of the Commandments. We may note also that Luther did not use the word \emph{missbrauchen}, misuse, as in his German Bible, but \emph{unnuetzlich fuehren}, take in vain. — The term “swear” has been misunderstood for so long a time that the correct view of it cannot gain ground without encountering considerable obstacles. As commonly understood, it means swearing in court. If this meaning is accepted, the word must be presumed either to forbid the taking of an oath altogether or the taking of a false oath. Luther, however, knows nothing of the first interpretation, declaring in his Large Catechism that taking an oath is “a very good work whereby God is glorified.” The second interpretation would be permissible; for the confessional
\quellseite{354}
booklets of the Middle Ages often contain the phrase “not to swear” as an equivalent of “not to swear falsely.” But the term “swear” occurs in these booklets also in the sense of thoughtless swearing in every-day life. Which of the two meanings Luther had in mind, will best be decided by his catechetical sermons of 1528 and the Large Catechism, drawn from these. Now, there is no doubt whatever that, in the Large Catechism particularly, all emphasis is placed upon “lying and deceiving by the name of God” and that false swearing, or perjury, is expressly mentioned as one form of lying by the name of God. It is also quite significant that the Latin translation of that passage in the Large Catechism which suggests most strongly the concise form of the Small Catechism (cf. Mueller, p. 397, § 62; Tr. p. 597), does not mention \emph{iurare} at all, but reads: \emph{illud blasphemando, execrando, maledicendo, incantando, ignominiose usurpamus}. The obvious conclusion is that Luther, in this connection, had in mind the thoughtless swearing in every-day life, a conclusion drawn also by Kantz, 1542; Meckhart, 1553; Trotzendorf, about 1556; the Joachimstal Catechism, 1574; Opitius, 1588; and others. Or, “swearing” might possibly be linked still more closely with “cursing” as a virtual synonym of that word so that the whole phrase would be practically equivalent to “using profane language.” The catechist who is guided by Luther’s intentions will treat of perjury under the head of lying and deceiving by the name of God; of the just, divinely warranted oath under the head of calling upon the name of God (cf. Large Catechism: Trotzendorf; Dressler, 1561; Hubernus, 1544; and others). He will thus avoid the pedagogic error of treating something that is commanded and classed as a good work, under the head of things forbidden (cf. Reu, Explanation etc., p. 36 f.)

In explaining the \emph{Third Commandment} the catechist will take special pains to set forth clearly the evangelical conception of the Sunday and to make plain that the Sunday is an institution arranged by the Church (cf. pp. 101 f.; 323 f.). For this reason he will not feel inclined to surrender the form prescribed by Luther, “Thou shalt sanctify the holy day,” and possibly to exchange it for “Remember the Sabbath day, to keep it holy.” He is aware, of course, that Luther did not originate the version which is found in his Catechism; and he also knows that even this version does not entirely correspond to the evangelical conception of the Sunday: on basis of Col. 2:16 we Evangelicals recognize no divine distinction between one day of the week and another. He also realizes that the adoption of the
\quellseite{355}
conventional English form, “Remember the Sabbath day, to keep it holy,” does not necessarily imply a denial of evangelical liberty; it is found, e.g., in all the Latin translations of the sixteenth century (also in the Book of Concord of 1580), and, through the influence of the Nuremberg Sermons for Children, also in several German reprints and enlarged editions (e.g., even in the \emph{Nuremberg Textbook for Children}, composed in 1628 and in use throughout the seventeenth century). Still, because it departs from the wording of the Old Testament, Luther’s version is better adapted to impress upon the minds of Lutherans their freedom from the literal phraseology of the Old Testament Decalogue. Formerly, perhaps, little depended upon the wording as such because Calvin as well as Luther rightly interpreted the words in question; but at this day it is well to pay more attention to a correct formulation. In the first place, the sections that went over to the Reformed Church (e.g., Anhalt, Lower Hesse, the Palatinate), felt called upon to “purge” or “reform” Luther’s explanation of the Decalogue by adding the words of the Old Testament, totally ignoring the New Testament freedom from the Jewish Decalogue (Reu i 1, p. 217; i 2, 2, p. 377; i 2, p. 278 ff.). In the second place, it is a fundamentally false attitude toward the Old Testament and the Law that is evidenced by the continued demands of the Reformed for a literal rendering of the Mosaic Decalogue. The catechist who would interpret the \emph{Third Commandment} in a really evangelical manner, will best equip himself for the task by a repeated perusal of the pertinent section in the Large Catechism. After emphasizing the correct attitude toward the Word and its proclamation and enlarging upon the blessing and importance of these, the catechist will find that his pupils do not despise the Church institution of Sunday; they will rather gladly meet together Sunday after Sunday, pray together, hear God’s Word, partake of the Holy Sacrament, and by that practise, manifest and edify themselves as a Christian congregation, and they become willing to lay aside all temporal labor that might prevent them from appearing at the common service, and quietly to meditate upon what has been heard and learned. Thus is made possible a Sunday rest that, quite apart from a purely human need of rest and the social aspect of the question, bears the impress of the gospel.—When the origin of Sunday as an institution of the Church is under discussion, it is customary to mention, alongside the fact of Sunday being the day of Christ’s resurrection, that Pentecost, too, fell upon that day (so also in our \emph{Explanation of the Catechism}, larger edition).
\quellseite{356}
But since this circumstance was no factor at all in the choice of just Sunday as the day of rest, that argument is irrelevant.

Although the several statements of the Fourth Commandment, grammatically, have reference to superiors as well as parents, they were chosen by Luther with a view to the latter rather than the former, a fact that the catechist dare not forget. Nevertheless, while enjoining upon his pupils the duties toward their parents, it will be incumbent upon him to flash the light of Scripture also upon the relation toward superiors, which is done best by referring to the Table of duties (\emph{Haustafel}). The duty of the catechist becomes so much more imperative as it is largely thought that, in this respect, the injunctions of the Bible have become more or less antiquated. Yet, however radical the change of social conditions, the duties of deference to authority and of faithful service stand unaffected. The idea—so classically expressed in the Large Catechism—that the parents are God’s own representatives and that all authority of superiors in school and state is grounded in the parental office thus defined, is of abiding moment. The better the catechist succeeds in arousing the conscience in this respect, the more his teaching will contribute toward preserving unimpaired the foundations of a healthy public life.

Rightly Kaftan points out that the \emph{Fifth Commandment} refers to the exhibition of the fear and love of God merely in the sphere of the neighbor’s bodily life; a teacher influenced by Luther’s example should not treat here of the commission of spiritual harm and murder. It is true that \emph{Leid tun} is not equivalent to German \emph{schaden} or our “harm” and that it actually denotes suffering in the sphere of the soul; nevertheless Luther refers in this connection to the neighbor’s body (we should fear and love God, and not hurt nor harm our neighbor in his body). Accordingly the subject under consideration is an “injury dealt to the body through harm done to the soul.” The right explanation would therefore be: “To embitter and shorten the life of the neighbor by envious looks, by spiteful words, or by hostile acts, as Esau and the sons of Jacob did.” When, however, Kaftan proposes to treat also of suicide under this head, he goes beyond Luther; quite likely also beyond the original intent of the Mosaic Law since in all the commandments, from the Fourth to the Tenth, the relation to our fellowmen is the dominant factor even though the term “neighbor” does not appear until the Eighth Commandment. Suicide had best be treated in connection with the Sixth Petition.
\quellseite{357}
In the \emph{Sixth Commandment}, the words, \emph{\emph{ein jeglicher} sein \emph{Gemahl} lieben und ehren}, are hardly to be construed as an additional main clause. The words are governed by the conjunction \emph{dass}, and \emph{ein jeglicher} refers back to the subject \emph{wir}: “and that each of us love and honor his spouse.” \emph{Gemahl}, “spouse”, as usual with Luther, is neuter although he probably had woman in mind (cf. \emph{ein jeglicher}). Present conditions and the duty to explain the text compel the catechist to discuss briefly the nature of marriage and to place strong emphasis upon its divine origin and its indissolubility. While the duties of husband and wife need merely be touched upon, the phrase “live chaste and pure etc.” requires a thorough as well as tactful treatment. Here Luther’s restriction to the positive element, of transcendent value pedagogically, should be the pattern for every teacher, especially when boys and girls are taught together. That which is forbidden “should be allowed to show itself only as the shadow, as it were, of that enjoined.” If there should be an occasion to enter into details, it should be done in private conversation with the individual. For the practical treatment of the subject much can be learned from Ziethe, v. Rohden, and Heydt, \emph{Die unterrichtliche Behandlung des 6. Gebots in der Schule. Drei gekrönte Preisschriften}, Berlin, 131894.

The \emph{Seventh Commandment} is to show us how the fear and love of God is to reveal itself in relation to the neighbor’s property. It is impossible to unfold Luther’s text without demonstrating the divine sanction of property and the unequal distribution of wealth as divinely intended. The teacher, likewise, must cast the light of truth upon the various forms of unfaithfulness or commercial dishonesty which his pupils may observe. The salient words serving this latter purpose are “nor get it by false wares or false dealing” (falsis mercibus aut impostura, in Sauermann’s translation; corruptis mercibus aut fraude aliqua, in the first; falsis mercibus, ullis malis artibus proximum decipiamus, in the third; fucosis mercibus aliisque actionibus fraudulentis nobis ea vendicemus, in the fourth Latin translation). “This applies to the indirect ways of stealing as it is done in all sorts of transactions, esp. in buying and selling, when the wares bought and sold are false, i.e., adulterated or counterfeit; or when the dealing is false i.e., when deception is practiced through misrepresentation, short weights and measures, etc., or when debts are not paid, either to the individual (loans, credit, keeping back wages), or to the company or state (fares, taxes, customs duties); or when trustees are faithless (embezzlement, breach of trust), or when advantage is taken of need (usury), or when
\quellseite{358}
no equivalent is given (gambling), or when false pretense is made (begging), and the like’ (Kaftan-Horine, p. 48); also practices of labor unions and trusts, insofar as they result in material loss to the neighbor, are entitled to attention should the occasion for a discussion of them arise (cf. our Explanation, p. 50). The catechist, however, should always bear in mind that he is to judge of all these things neither from the legal nor from the economical, but from the religious standpoint so that, also here, not natural, but Christian, Evangelical ethics are inculcated. Here, too, Luther has placed in the foreground the fear and love of God as the source of power and the motive of action. He who fears and loves God will recognize the divine order, also in regard to property, and beware of displacing the divinely imposed bounds and barriers; for he knows that he has in his heavenly Father the treasure of treasures and a helper in poverty and distress, and that earthly property has never permanently satisfied the heart. Only the lover of God can help his neighbor to improve and protect his property and business: he has learned from the heavenly Father that it is more blessed to give than to receive. The term \emph{Nahrung} here does not signify food, or provisions, as in the First Article, but is used in a wider sense, signifying the means of livelihood, the source of income, trade, business (cf. Gen. 46:33 f. in the German version), and, in general, income, earnings, property (\emph{conditio}, second Latin translation). The positive part of the explanation does not contain three statements: help, improve, protect; but only one: we should help to improve and protect our neighbor’s property and business (Sauromannus: \emph{sed demus operam ut illius opes conserventur et eius conditio melior reddatur}; H. Magdeburg: \emph{sed ad augendas facultates eius et conservandam rem familiarem prompti simus}).

According to the \emph{Eighth Commandment} we should exhibit our fear and love of God by exemplifying both truth and love in all that we say about our neighbor. Well it is said by Th. Harnack in his explanation of the Catechism (in \emph{Katechetik} ii, Erlangen, 1882): “Truth and love must always be found together. Truth protects love from the fault of good-natured indulgence and weakness which arises from indifference to the truth of the matter. Love, conversely, protects truth from the severity of injustice arising from indifference to the person. For this reason truth is the light and salt of love, and love is the fire and warmth of truth”; compare in connection with this and the other Commandments, in addition to Luther’s Large Catechism, the Nuremberg Sermons for Children (Reu, i 1, pp. 462-564), a work always
\quellseite{359}
instructive, though occasionally going beyond Luther. \emph{Beluegen}, “belie”, according to Luther’s usage, does not by any means signify to lie to anyone, but to lie about him, to spread lies about him; “falsely” is added not only to add force to the conception of falsehood conveyed by the verb “belie”, but also to disclose the motive back of such falsehood,—to speak as a deceiver, purposely and craftily to utter a falsehood in order to harm the neighbor. “To betray” means to imperil some one through disclosure of what he desires to keep secret; it is equivalent to Latin \emph{prodere}. This being its meaning, the example of Judas, frequently adduced, is hardly suitable. The motive of Judas was quite different; nor did Jesus desire the place where He kept Himself to remain concealed. \emph{Afterreden oder boesen Leumund machen}, Luther writes and thus connects these two verbs as closely \emph{luegen and truegen} in the Second, or \emph{abdringen and abwendig machen} in the Tenth Commandment. The slandering (\emph{boesen Leumund machen}) is done by means of backbiting (afterreden); by speaking evil about the neighbor behind his back, the coward brings him into bad repute. Hiob Magdeburg combines both notions into one in his splendid translation: “\emph{Timendus et amandus Deus est ne de proximo dicamus mendacium aut invite illum deferamus neve existimationem eius violemus lingua futili, sed ut illum excusemus, bene de illo loquamur et in meliorem partem interpretemur omnia}”. Sauermann retains also the adverb “falsely”: “\emph{ne proximum falsis mendaciis involvamus}.” The German word \emph{Leumund} is by no means an abbreviation for der Leute Mund; its Old High German equivalent is \emph{hliumunt}, from the same root as the Gothic \emph{hliuma}—ear, hearing; so \emph{Leumund} is—\emph{fama}, Engl. fame, and \emph{boesen Leumund machen}—defame. The Joachimstal Catechism (1574; probably based on Mathesius) says: “ ‘Falsely belie’ means to raise false accusations against one before the authorities and other people, to impute to others the abuse of which oneself is guilty, like Gehazi. ‘To betray’ means to divulge everything, to peddle around what we have heard, and to be mongers of scandal; likewise, without necessity to reveal secrets, as the Ziphites did in the case of David. ‘To Backbite’, to make known behind their backs what people do and say; to censure and put the worst construction upon it, as did Doeg. ‘To slander’ means by cutting remarks to hack his honor to pieces as meat is hacked to pieces upon the butcher’s table with the result that thus evil suspicion fastens to him and he is robbed of his good name as was done by Absalom in the case of his father David.” Tetelbach, 1568: “ ‘To excuse’ means to defend the neighbor against evil mouths and false tongues; ‘to speak well
\quellseite{360}
of him’, to evince a kindly disposition toward him wherever he is the subject of conversation; ‘to put the best construction on everything’, to cover up the frailty and fall of the neighbor and to excuse and explain them as best we can to those who speak ill of them.” Siber, 1575: “ut ad calumniatores excusemus, obtrectantes refutemus, et commode sentientes de illo in melius omnia, interpretando accipiamus.” Tetelbach, who understood Luther, takes pain, in connection with every commandment, to point to the root of true fulfilment—the fear and love of God; here he asks the question, Who transgresses this commandment?, and his answer is: Those who do not fear God, but follow the devil, the father of all liars and slanderers. Afterward he asks the other question, Who fulfills this commandment? Answer: Those who love God; for the love of God works in us truthfulness, sincerity, and a helpful tongue (Reu i, 1, p. 681).

As to the Ninth and the Tenth Commandments, Luther had previously been in the habit of treating them as one, even in the Large Catechism. If now, in the Small Catechism, he treats them separately, the reason is, of course, not that in the meantime it had dawned upon him that the verb used in the Ninth Commandment forbids \emph{concupiscentia inheridata}, and that used in the Tenth, \emph{concupiscentia actualis} (cf. Deut. 5:18, \emph{chamad} and \emph{awah}). In previous expositions he had brought both Commandments in connection with original sin, but here he says nothing whatever on that subject. Nor did he keep them apart because he realized that the object in the Ninth Commandment is in the singular while that of the Tenth Commandment is in the plural number; there is no evidence whatever that he had this difference in mind, and his contraction of two of the objects named in the Bible—ox and ass—into one—cattle—is an argument against this view. Luther’s motive was probably concern for the young that had to learn these Commandments: they would have been bewildered at finding no separate explanation for each Commandment as long as they were held to number them separately, in accordance with the traditional practise. The catechist who does not aim to supplement Luther’s Catechism, but merely to unfold it, will, for this reason, hardly take the pains to find a difference between the Ninth and the Tenth Commandment, especially since, until today, no real difference has irrefutably been shown to exist. But if he feels the need to explain why the Tenth Commandment is numbered separately, the method of the Nuremberg Sermons for Children commends itself as most practicable, provided we shall not bring in the difference between inherited
\quellseite{361}
and actual lust. There we read: “Having learned in the previous commandment that we should not desire to dispossess the neighbor of his house, position, or calling, we might come to the conclusion that while we should not desire his house as a whole, to desire a part of it, such as a good servant or ox, were no sin and that we might well try to accomplish our desire. Lest we should entertain such a thought, God forbids that, too, and says: ‘Thou shalt not covet thy neighbor’s wife, nor his man-servant, nor his maid-servant, nor his ox, nor his ass, or anything that is thy neighbor’s.’ Otherwise the thought might come to us: Oh, even though I take this or that from my neighbor, it will not harm him, for he has enough of such or similar possessions left. But the Lord does not want this; hence he summarily forbids us to covet anything whatever possessed by the neighbor.” In this connection color may be given to the explanation of the Ninth Commandment by the reflection that the house or piece of ground inherited by the Israelite had an added value in that it was to him a guarantee of his share in the promised land (cf. the laws bearing upon the year of restoration; see also Kirch. Zeitschrift, 1917, p. 372 ff.).

More important than the question as to the difference between the Ninth and the Tenth Commandment is that as to the progress made by these two upon the preceding ones, especially the Sixth and the Seventh. According to the words of Scripture, the element of progress is indicated by the verb. In the previous Commandments, the wicked word and deed are forbidden; here God forbids the wicked desire; but desire, whether we take it as inherited lust or as actual lust, is still a matter of the inner world of the heart. That these two Commandments were, in consequence, of special importance for the Jews who fondly thought that they had kept the Law if only they did not transgress it in word and work, is self-evident. More difficult it was for Luther to find something new in the last two Commandments since he was of the opinion that every Commandment is to be applied to the heart. The element of progress, accordingly, could be found only in this that the Ninth and Tenth Commandments contain an injunction against inherited lust. He writes, e.g., in Decem Præcepta Wittenbergensi prædicata populo (1518): “Videtur . . . duobus præceptis prohiberi ipse fomes et invincibilis cupiditas, ipsa, inquam, radix malarum cogitationum, ut sc. 6. et 7. intelligantur prohibiti consensus cordis et signum membrorum, sc. verbum oris et opus corporis mali, hic vero etiam ipsi primi motus una cum fomite, qui est origo illorum”; he refers to Rom. 7: 23 f. Taking
\quellseite{362}
this view, it is indeed possible for him to say that the fulfillment of the other Commandments is possible while these two are impossible of fulfillment for anyone, no matter how holy he is, and that in their light one becomes aware of his sinfulness. Through them man learns to despair of himself and his merits, and to trust in the mercy of God. So Rom. 7:7 and Luke 11:39 would be the key to the matter. While, in the Large Catechism, he holds fast to the conclusion that the element of progress is signalized by the verb, the word “covet” is to him now no longer an expression of original sin, but of actual lust, of lust to which man’s will has yielded (although he admits incidentally “that the scope reaches somewhat farther and higher”). He writes, “God has added these two Commandments in order that it be esteemed as sin and forbidden to desire or in any way to aim at getting our neighbor’s wife or possession.” It is required that no one “should think or purpose to obtain for himself what belongs to another, such as his wife, servants, house and estate . . . even with a show of right or by a subterfuge, yet with injury to his neighbor” (\emph{Weim. Ed. xxx 1, p. 175; Tr. p. 665}). Instead of original sin, he now recognizes avarice and envy as the root of all sinful conduct against the neighbor. “Therefore we allow these commandments to remain in their ordinary meaning, that it is commanded, first, that we do not desire our neighbor’s damage, nor even assist, nor give occasion for it, but gladly wish and leave him what he has, and besides, advance and preserve for him what may be for his profit and service, as we should wish to be treated. Thus these Commandments are especially directed against envy and miserable avarice, God wishing to remove all causes and sources whence arises everything by which we do injury to our neighbor, and therefore He expresses it in plain words: Thou shalt not covet, etc.” (\emph{Weim. Ed. xxx 1, p. 178; Tr. p. 669}). The impression is received that Luther, in consideration for the common man and children, has confined himself to such “ordinary meaning.” Nor is it likely that he intends to turn aside from this meaning in the concluding words: “God would especially have the heart pure, although we shall never attain to that as long as we live here; so that this commandment will remain, like all the rest, one that will constantly accuse us and show how godly we are in the sight of God!” (ibid.) What he means to say, is in all probability this that we shall not even succeed in attaining to such a degree of godliness that we never yield to the evil desire and never begrudge our neighbor his possessions. In the Small
\quellseite{363}
Catechism he takes his cue altogether from his statements in the Larger Catechism. In those of the preceding Commandments which have reference to our conduct toward the neighbor, he has as a matter of fact (Kaftan to the contrary notwithstanding) merely spoken of the words and works in which our fear and love of God is to find expression, but not, as yet, of the control that such fear and love is to exert over our thoughts in regard to the possessions of the neighbor. This, however, he now supplements in the Ninth and Tenth Commandments, where he explains “covet” as denoting “to seek to gain.” To be sure, he does not stop there, but at once proceeds to enumerate all the evil practices to which one is but too likely to resort when he seeks to gain the neighbor’s possessions; but by using in his explanation the term “to seek to gain,” the verb, and herewith the element of progress in these commandments, has not remained entirely unexpressed. The catechist that desires to remain in Luther’s footsteps will make the concept “covet” his starting-point, and explain: “Covet means to envy the neighbor his possession and to be discontented until it has come into one’s possession.” Then he shows how such coveting leads to the acts mentioned by Luther in his explanation. While he thus would not find original sin forbidden here directly he would come to speak of it as he develops the subject; for he would surely have occasion to make clear to the class how it comes that covetous thoughts will arise in us again and again (cf. our \emph{Explanation}, p. 55).

Just in view of the persistence with which covetous thoughts arise in us, fear and love are to be manifested in not only opposing them, but in doing one’s best to help the neighbor preserve his possessions. We read in the \emph{Nuremberg Sermons for Children}: “Since we are aware of the fact, my dear children, that coveting is a sin, let us guard against it, by the grace and help of God, as much as we can. But we should take special pains not to yield and conform to any evil desire . . . God has created everything and is its true Lord; therefore He gives it to whom He pleases. Has He given it to your neighbor, let him be welcome to it, and remember: if God had wanted me to have it, or if it had been of any benefit to me, He might have given it to me, too” (Reu i 1, p. 503). Tetelbach: “He who loves God does not begrudge his neighbor any blessings but is rather of help to him in keeping them” (Reu i 1, p. 683).

In regard to “house”, the Nuremberg Sermons say: “You must not merely understand the building by the little word house—the structure that the people live in, but the whole household and everything
\quellseite{364}
that pertains to it. For, in the Old Testament, the cities and villages were distributed among the Jews according to tribes, and if any one did not belong to the particular tribe occupying a certain place, it was impossible for him to obtain a house there. Therefore we may understand by the word house as much as tribe. This, therefore, is the meaning: If our neighbor belongs to a superior tribe or family; if he has inherited much property; if he is the owner of a fine household, possessing citizenship, honors and dignities, an honorable office, and everything pertaining to a householder, we should not covet his house, i.e., we should not have the least desire that he should lose it, so that we might take his place“. (Reu i, 1 p. 501 f.). Tetelbach: “By inheritance any kind of real estate is meant, such as fields, meadows, gardens, vineyards, timber, ponds, country-seats, and possessions of that sort. House means the dwelling of the neighbor and everything outside and inside that pertains to it, for instance, furniture. To seek craftily implies every kind of intrigue and wicked practise, bringing the neighbor to grief, depriving him of his income, damaging his property, or compelling him to sell out and to assign his possessions to another. Show of right means to lay claim to property by fraudulent titles and false pretenses; to outrage the law and to precipitate legal action; to bribe unscrupulous people; to advocate the side known to be wrong; to pervert the law through financial maneuvers in order to secure a favorable decision and thus to obtain the neighbor’s property”.—The second half of the Ninth Commandment should be construed in this manner: “but help and serve him in keeping it”. Compare the first Latin translation, Sed potius \emph{iuuare eum ut suas fortunas retineat integras}; the second, Sed \emph{detur opera sedulo ut ista proximo diligenter custodiantur}; the fourth, Sed \emph{ad conservanda ea ut prompti et parati simus}. “Estrange” means to alienate by allurements or promises; the German word \emph{abspannen} is by no means identical with \emph{abspannen}=to unhitch; its root is either \emph{spannen}=persuade, lure, entice, deceive (cf. High German \emph{Gespenst}), or \emph{spenen}=to wean (cf. \emph{Spanferkel}—a little pig recently weaned). Luther preferably connects the word with wife (cf. Weim. Ed. xxx 1, pp. 9. 42. 84). Tetelbach, while mistaken in regard to the derivation, gives the correct explanation: “With smooth words and promises to persuade them to neglect their calling and duty and to cease doing right, as when a horse is unhitched or a crossbow unstrung.” “To force away” Tetelbach explains as follows: “By threats, force, and undue pressure to urge them until they leave the neighbor”. “To entice away” means, according to the same author, “to steal one’s
\quellseite{365}
heart and favor by gifts, matchmaking, and other unscrupulous measures so that he turns against his environment and wishes for another.” The first Latin translation has “\emph{ad nos arte aliqua transferre, non vel vi vel dolo aliquo eos ab aliis abalienare}”; the second, “\emph{abalienemus aut abstrahamus}”; the fourth, “\emph{ne ulla ratione abalienemus et ad nos alliciamus uxorem proximi, servos, aut quicquam ex ipsius familia}.” Kaftan is right when he says: “All three expressions have substantially the same significance; ‘to entice away’ is the general term, ‘to estrange’ refers to the person enticed away and denotes that this is accomplished by allurements; ‘to force away’ refers to the owner and denotes the use of persistent urging.” That Achelis is mistaken when he applies the term ‘estrange’ to animals, ‘force away’ to servants, and ‘entice away’ to the wife (in Der Dekalog als katechetisches Lehrstueck, p. 53), requires no argument.

In regard to the function assigned by Luther to the Conclusion, Albrecht (in Weim. Ed. xxx 1, p. 361) recapitulates the matter bearing on this subject as follows: “In deference to tradition, Luther left out of his text the prohibition of images together with the threat and the promise. That he accounted these words as parts of the First Commandment, is evident from his comments on Exodus (1525, cf. Weim. Ed. xvi, pp. 436 ff.; 445, 23) where we find the significant word: ‘The Promise and Threat Contained in the First Commandment.’ Aside, however, from this significant statement, Luther, in his earlier catechetical expositions, failed to take note of this addition, although he maintained from the outset the fundamental, controlling importance of the First Commandment. Not until the second and third series of Catechetical Sermons of 1528 (W. E. xxx 1, pp. 43. 85), did he take the addition into consideration. Here he emphasized the fact that the addition was virtually part of each commandment (ibid., p. 43), and that he had appended it at the end of the commandments for the purpose of ‘joining them into a wreath, the last to the first.’ The threat with its power of arousing fear and the promise with its power of arousing love are to him the ‘cord with which the wreath is tied together’ (ibid., p. 85). Quite in harmony with these statements is the exposition in the Large Catechism where the conclusion receives attention both in connection with the First Commandment and after the last” (ibid., pp. 136 ff., 179 ff., Tr. pp. 589 ff., 669 ff.). That this does not justify the inference that the “fear” of the First and the succeeding commandments is to be understood in the same sense as in the conclusion, has already been dwelt
\quellseite{366}
upon (p. 327 f.; compare also A. Hardeland, \emph{Der Begriff der Gottesfurcht in Luthers Katechismus}, Guetersloh, 1914, against J. Meyer, \emph{Fuerchten, lieben und vertrauen in Luthers Kleinem Katechismus}, in \emph{Neue Kirchl. Zeitschrift}, 1913, pp. 793 ff.). In regard to the position of the conclusion see also p. 101. The catechist who favors Luther’s arrangement of appending the threat and the promise to the Commandments as a whole will certainly let his treatment of the Commandments culminate in an exposition of this threat and this promise as two powerful incentives to obedience. To raise the question at this point in regard to “the use of the Law”, is to such a catechist an incomprehensible procedure. Has he not in all preceding catecheses tried to direct his pupils’ eyes to the Law as the norm of life, which certainly is a use of the Law? No one would ever have thought of introducing such a discussion at this point but for the serious misconception that dogmatical topics must be interpolated in the Catechism. Again, the catechist following in Luther’s steps will not think that he must pause to show the children their sin in the light of the Law, true as it is that the Ten Commandments can serve as the mirror for sin, and should serve for it at the proper time, namely in connection with confession. Still less will he now take the opportunity to convince them of their inability to keep the commandments. Nobody would have been led to these two pedagogic follies had it not been for the desire to systematize the Catechism by linking the two Chief Parts together. On the contrary, just as it was the catechist’s aim to show his pupils in connection with each Commandment how child-like fear and love of the Father should impel them to the fulfillment of the divine will, even so he reminds them, for the same reason, of the wrath and jealousy divine which is to be visited upon the transgressors of the Law, and of all the blessings in which the obedient may rejoice, in order to help them by all means to learn obedience to the divine will. A treatment of the Commandments which conforms to Luther’s pattern, dare not sound as final note the confession: “I am a sinner, groaning under the curse”, nor the conviction: “I am too feeble to fulfill the Commandments of God”. Rather should it lead to the resolution: “As the child of my heavenly Father I mean to fulfill the will of my God and Father, and to fulfill it cheerfully” (it is just this voluntariness and cheerfulness of an obedience prompted by fear and love which Luther emphasized again and again), and to the prayer: “Help me, my God and Father, better and more perfectly to fulfill Thy will!”

The text of the Bible contrasts those that love God with those
\quellseite{367}
that hate Him. It speaks of hating because, according to its very nature, sin is rebellion and hatred against God (Rom. 8:7), a fact quite in evidence in the case of those who wilfully and defiantly disregard the commandments. “Unto the third and fourth generation of them that hate Me” and “showing mercy unto thousands of them that love Me” are phrases bringing out the energy, the earnestness, and holy zeal of God, both when He inflicts penalties and when He scatters blessings. The thought of them is to add strength to the child-like fear and love in our hearts toward the heavenly Father as an incentive to obedience. Both punishment and blessing bear primarily upon the bodily life; the fundamental principle is brought out that sin is followed by evil, obedience by blessing. In connection with transgression we are to think of every kind of misfortune, failure, trouble, sickness, etc.; and in connection with obedience, of blessing in body in that God preserves our health and strength, of blessings in soul in that He rewards obedience with peace and joy, of blessings in property in that He prospers our work, of blessings in honor in that He gives us a good reputation among all godly men. But because these results are not invariably the instantaneous fruit of obedience, Luther, once more, rightly lays stress upon trust. If God permits us to languish in misery in spite of our obedience, it behooves us to trust Him for finally changing suffering into bliss, even though we should have to wait till eternity’s day. The words employed are broad enough to embrace this last thought. Well it is said in the Nuremberg Sermons for Children where the Conclusion is brought in connection with every one of the Commandments, and each sermon ends with the following words:

\begin{quote}
“Thus, my dear little children, you have the correct, simple, everyday meaning of this Commandment. You should lay it diligently to heart and fear God lest you transgress this commandment; for what He commands is good and holy, and what He forbids is wrong, sinful, and shameful. It is the determined will of God that we should keep His commandments and not despise them; for He says: ‘I, the Lord, thy God, am a jealous God’. . . . Therefore, my dear little children, you should fear God and diligently keep His commandments, imploring His grace and help so that you may be enabled to do this. For the fear of the Lord is the beginning of all wisdom and makes us pious and good, pleasing in the sight of God and of service to other people; and we shall enjoy peace and quiet, honor, prosperity, and happy days. If we thus continue in obedience to God and in the right faith to the end, God will add eternal life as the last gift. May He grant
\quellseite{368}
this to us all. Amen.”
\end{quote}

This is the final note to be sounded in an exposition of the Commandments. Should anyone think that we should thus turn into the Roman channel, he merely reveals his failure to understand either the Scripture or Luther. Albrecht says aptly: “There is no trace here of pleading of merit; trust clings to a promise of grace, not of a reward” (W. E. xxx 1, p. 362). The believer, justified by faith, who, if all is well with him, endeavors to fulfill the will of God by inner necessity, finds an additional motive and incentive in the “grace and every blessing” promised of God upon the sole condition that we obey.

The catechist will much better succeed in impressing upon the heart of his pupil the thought that the regulations in the Father’s house, as found in the First Chief Part, are to be faithfully and scrupulously kept, if he sums up each Commandment by reaching into the life of Jesus, the perfect fulfillment of the Law. The fulfillment of this or that single requirement is typified by one person of sacred history or another; but of the fulfillment of every requirement we have in Jesus the most winsome and attractive type: all rays of obedience meet in Him; not one feature is lacking. Such treatment as this will arouse in the catechumen the conviction that the object to be realized is to follow Christ; that it is His image that is to shine forth from our life. Thus the New Testament character of the First Chief Part stands duly revealed. Leopold Schulze has laid great stress upon this thought in \emph{Katechetische Bausteine} (Magdeburg, ⁵1891); in our own \emph{Explanation of the Catechism} it has been carried out consistently.

The needlesness of a transition to the \emph{Second Chief Part}, and Luther’s attitude in the matter, have been discussed above (pp. 111 f.; 319). Here we only add what Albrecht says on the subject (Weim. Ed. xxx 1, p. 362): “In the Small Catechism the \emph{Second Chief Part} is adjoined to the First without any connecting thought. This is the case also in the third series of Catechetical Sermons of 1528 where we read in the sermon of December tenth: ‘You have heard the first part of Christian doctrine, namely, the Ten Commandments. I have admonished you diligently that you should exhort the family to learn that part literally by heart so that they might obey God accordingly, and you as their masters, and yourselves likewise might obey God. For if you will instruct and urge the family, there will be progress; and there has never been a doctor who has failed to become more learned by study. These Ten Commandments having been dealt with, we will now proceed to the next part’ (xxx 1, p. 86). At other places
\quellseite{369}
connection is made by harking back to the First Commandment, or by asking the question, \emph{Qualis est deus tuus qui tulit tibi decem praecepta?} (xxx 1, pp. 9. 10. 183), or by reviewing the connection given in the short form [above pp. 85, 111 f., 319]: those incapable of fulfilling the commandments [for the disease of sin does not here come into consideration as guilt, as is still maintained by many, but rather as moral impotence] will receive help and strength from God through faith (xxx 1, p. 43 f.; xi, p. 48). This thought though expressed less forcibly, returns also in the Large Catechism (xxx 1, p. 184).” One who explains the Introduction, the First Commandment, and the Conclusion according to Luther’s plan, has learned that the First Chief Part presents all of Christianity, insofar as it is a divine requirement for the baptized Christian, and he will not only feel no need for the interpolation of transition material but will find it absolutely incongruous.

In the \emph{First Article} Luther places all possible emphasis upon the word Creator or “Maker.” Not only does he give this article the heading “Of Creation”, but, in the edition of 1531, revised and corrected by him (cf. p. 95), he had the printer capitalize every letter of the word Schoepffer (just as the word Herr in the Second Article). He made it the keyword in his explanation. Previously his mode of operation had been different. In the Short Form he did not emphasize the objects of faith, but the verb—the concept of faith, although mention is made of the former (cf. p. 90). In the sermons of 1523 he lays more stress upon the objects of faith, especially on God as the “almighty Father” (Weim. Ed. xi, p. 49); similarly in the rather brief allusions in the first two series of sermons of 1528. In the third series, however, Luther confines himself altogether to the word “Maker”: “The children and the common people are to learn the matter in short, plain, and simple terms as follows: There are three articles in the faith, (1) of the Father, (2) of the Son, (3) of the Holy Spirit. What do you believe of the Father? Answer: He is the Creator. What of the Son? He is the Redeemer. What of the Holy Spirit? He is the Sanctifier. For the benefit of the learned the article might be divided into as many parts as there are words; but now I would teach the young and the common people that God the Father is the creator of heaven and earth. Now, what is the meaning of that? Just this, that I am to believe that I am God’s creature; that He has given me body, soul, sound eyes, reason, property, wife, children, fields, meadows, hogs and cattle. Then that He has given the four elements, etc. This article, accordingly, teaches that you do not have your life of yourself, not
\quellseite{370}
even a hair. All things in existence are comprehended in the little word Creator” (Weim. Ed. xxx 1, p. 87). Such is the tenor also in both the Large and Small Catechisms although here, in the phrase “fatherly goodness and mercy,” the word “almighty Father” comes into its own. Whether this thought served as basis for the following statements “clothing and shoes . . . . . from all evil” is doubtful.—That Luther used the word “almighty” as an adjective, connecting it with “Father,” not with “creator,” which would be a tautology, hardly requires any proof. In the edition of 1531, so carefully revised by him, he puts a comma after “almighty”; also in the first edition of the Large Catechism (W. E. xxx 1, p. 130); then, after repeatedly omitting all punctuation, again in the reprint of the third edition (1540) and all subsequent prints; with unmistakable clearness it is indicated in that part of the exposition where we read: “But to speak more of this belongs into the other two parts of this article where we say ‘Father almighty’” (W. E. xxx 1, p. 184.: Tr. p. 681, § 17, 18). Capito (1528) and Brenz (1529) write without hesitation “almighty Father”. See M. Reu, \emph{Die Konstruktion des ersten Artikels} in Kirchliche Zeitschrift, 1904; Cohrs, \emph{Katechismusversuche vor Luthers Enchiridion} iv, pp. 295 f. That the right connection was retained for a long time after Luther, can be abundantly demonstrated from our sourcebook, \emph{Quellen zur Geschichte} etc.; why Luther did not simply write “almighty Father”, transposing the adjective as Capito and Brenz did, we have seen (p. 100).

Another problem is closely connected with, and to some extent solved by, the problem just considered. We refer to the question whether God comes into consideration primarily as “Father of Jesus Christ” or as “our Father.” Now, God is unquestionably considered here, in some way, also as the Father of Jesus Christ; this appears from the origin of the Creed as a baptismal confession, and from any number of New Testament passages which show that God is our Father only through Jesus Christ, His incarnate Son, coequal and coeternal with God the Father. But that this relation is to be chiefly emphasized at this place, as is conventionally done, or that it is the sole fact to be considered, is totally erroneous and most certainly not in harmony with Luther’s views. He writes in the Short Form: “Because He is God, He can do the thing that is best for me, and knows what that thing is. Because He is Father, He wills to do what is best for me, and to do it with all His heart. Because I do not doubt, but put my trust in Him, I am assuredly His child, His servant and His heir forever, and as
\quellseite{371}
I believe, so will it be done unto me.” In the second series of his Catechetical Sermons of 1528 we read: “\emph{Pater hat sich ausgeschuettet cum omnibus creaturis quas creavit}.” In the Large Catechism the same thought occurs in this form: “And He does all this out of pure love and goodness without our merit, as a benevolent father who cares for us that no evil might befall us” (Tr. p. 681). It is found once more in the hymn of faith: “\emph{Wir glauben all an einen Gott, Schoepfer Himmels und der Erden, Der sich zum Vater geben hat, Dass wir seine Kinder werden}.” In his \emph{Hauspredigt von den Artikeln des Glaubens, in Schmalkalden gehalten} (1537) he says: “Moreover, He not only teaches us what and whence we are, but also where we belong. This is indicated by the word “Father”, which means that God lays claim to fatherhood as well as to omnipotence. The animals cannot call Him Father; but we can call Him Father and ourselves His children. With that word He expresses His purpose in regard to us. Previously we had learned whence we are and what honor and glory are ours; now the question arises as to our destiny. It is this that you shall be children, and I Father. I have not only created you and seen to your preservation: I want you for my children and appoint you as my heirs so that you shall not be thrust out of the house like other creatures—oxen, cattle, sheep, etc., which, without exception, die or are devoured; but in addition to being My creatures, you shall be also My children for eternity. This is our prayer and confession when we recite the creed: I believe in God the Father” (W. E. xlv, p. 16). This was the sense in which, before the appearance of the Small Catechism, Brenz, Althamer, Lachmann, and others interpreted the word “Father.” This was its interpretation also after its appearance. Matthaeus Zell (1536) writes: “What does he mean when he says: Father? This is the most comforting thought that we may entertain concerning God. Why? In this one word the truth is expressed that the high divine Majesty loves us poor worms as heartily as a father, and not as a stepfather; that He desires us for children and for heirs of an eternal blessed life—a thought than which there is no higher and surer comfort in the extremity of death and in any other ills for anyone who, in true faith, speaks and confesses that word . . . . What do we mean when we confess that God is almighty? First of all, the man that really lays hold of that truth finds it a comfort exceedingly great. How is that? Not until we realize that He is almighty, have we the sure comfort that no one can prevent Him from doing as He pleases, unlike human parents who are often prevented from carrying out their will. But
\quellseite{372}
God, being almighty, can be prevented by no one: no matter when His help is required or by whom it is required, no one shall be too strong and powerful to be in His way” (Reu i 1, p. 109). The \emph{Nuremberg Sermons for Children} are less clear on the subject; but they give a splendid explanation of the term “almighty”: “Nobody is so sick but He can heal him; nobody so poor but He can make him rich; nobody so despised but He can bring him to honor; nobody so great a sinner but He can make him godly; nobody so bereft of faith but He can make him a believer; nothing so incredible but He can perform it if He will” (Reu i 1, p. 69). Butzer is clearer (1537): “I believe that God, the Father of our Lord Jesus Christ and the first person of the Godhead, desires to be my Father also. How is He your Father? In that He imparts to me His nature and heritage. What is God’s nature? All righteousness, godliness, and love. What is His heritage? Eternal, blessed life” (Reu i 1, p. 69). Huber (1544): “Therefore I believe and confess that Almighty God, the only, eternal God, is the almighty Father of myself and all the faithful” (Reu i 1, p. 792). Meckhart (1553): “Why do you call God your Father? Because He, through grace, has elected me as His child, promised me fatherly faithfulness and every blessing, and made me an heir of all His possessions” (Reu i 1, p. 824). But even those who take the word in the trinitarian sense, do not overlook the other aspect. A case in point we find in Kantz (1542): “Why do you call God Father? Because He has a Son of the same essence as His own, born from eternity, who is God equal with the Father . . . And through the same beloved Son He has elected as children all those who believe in Him so that they may inherit His kingdom with Him; John 1:12; Rom. 8:17. Such is the meaning of the word Father which is an exceedingly gracious, comforting word, assuring us that God is our faithful Father, and will be such forever; Isa. 49:15” (Reu i 1, p. 608). Or Tetelbach (1568): “Why do you call Him Father? He is such; for He has brought forth from His divine nature and essence a Son—our Lord Jesus Christ through whom He has also become my dear Father and the Father of all believers in the Lord Jesus Christ; and He is willing to deal with us as does a father with his children” (Reu i 1, p. 658).

Understanding the First Article thus we can comprehend why Luther, in 1523, could assign first rank to the First Article: “Hic primus articulus, quod credo Deum esse patrem meum, qui articulus summus est . . . Quicquid est in tota Scriptura, huc referri potest. Si deus pater meus est, liber sum ab omnibus, inferis, morte, mundo etc.” (Weim.
\quellseite{373}
Ed. xi, pp. 50 f.); and in that very connection he makes the words “almighty Father” the nucleus of his whole explanation. It is quite plain that Luther did not intend to cast aside Christ’s death and resurrection and the atonement conditioned by them; Luther says often enough that God is our Father only through Christ. And it is only when this is borne in mind that Harnack’s statement is acceptable: “The First Article is for Luther an epitome of Christianity” (\emph{Dogmengeschichte} iii, p. 711, note 3). From this standpoint another matter becomes clear, namely, that it is a thorough misunderstanding of the First Article when the view is still entertained as though it were nothing else but general faith in providence, made so much of in the era of rationalism, which is here confessed, and as though the First Article represented a kind of preparatory stage, a sub-Christian creed which rationalist and Lutheran Christian could very well join in confessing. Nothing could be more erroneous. That a faith in providence is here spoken of, is true; but it is the faith in providence as entertained by the Christian—a faith based upon his faith in redemption, yes, a faith which, far from being inferior, is nothing less than an expression of it. It is God, the almighty Father and Creator, of whom I am here speaking, who has become my Father through no one save Christ, and only through His redemption. Likewise, it is the baptized Christian that here confesses his faith, already a child of God through baptism, and now looking up to Him as Father and finding comfort in His care and protection. Instead of being something sub-Christian, the assurance of faith here manifest is specifically Christian and radically different in principle from the natural man’s faith in providence. It finds its true parallel in that stanza of Paul Gerhardt’s: “With gladness I confess it, What long my heart has known: That God, most high and blessed, To me His grace has shown; And that, whate’er betide me, He by my side shall stand, Through storm and wave to guide me, And every ill at hand.” Only thus can we justify the fullness of assurance of Luther’s words: “He richly and daily provides me with all that I need for this body and life, protects me against all danger and guards and keeps me from all evil.” The more the uniqueness of the Catechism is preserved also in this respect, the more clearly the gospel will stand disclosed.

In regard to the explanation of the First Article, a difficulty arises from a mooted point of construction. So much is clear that Luther’s explanation is readily divided into three parts: (1) What God my Maker has done, and still does, for me; (2) Why God has done and
\quellseite{374}
still does this; (3) What I owe him therefore. But how to construe, in the first part, the words \emph{dazu Kleider und Schuh . . . . . reichlich und taeglich versorget} is the question. Frequently it is advocated to subordinate the whole section \emph{dazu Kleider und Schuh . . . . . bekleidet und bewahrt} to the words \emph{und noch erhaelt}, and still preserves. Divine preservation takes place when God supplies everything which I need for the preservation of life and when he safeguards me against everything that might cause me injury. Such, it seems, was the construction preferred by H. Opitz (1583) who answers the question, “Whereby does God preserve our human body and life?” in this manner: “By clothing, shoes, . . . cattle and other goods, together with all that is needed to support this body and life” (Reu i 1, p. 532). With him agrees H. Magdeburg, so far as his Greek translation is concerned (the Latin one is different: v.i.): “\emph{καὶ νῦν ἔτι τρέφει, ἐπιμελούμενός μοι τῶν ἐσθήτων, τῶν ὑποδημάτων . . . . . οἰκίας, πρὸς δὲ τούτοις τῆς γυναικός . . . . . πρόνοιαν ἔχων, ὅπως ἂν ταῦτα ἱκανῶς καθʼ ἡμέραν ἑκάστην ὑπάρχῃ}.” The result of this construction was that the word \emph{dazu} (= also) was no longer taken in the sense of “further”, but connected with the preceding “still preserves to me”, and thus understood to mean “for this purpose”, namely for the purpose of preservation. After the words “and all my goods” a semicolon was placed, the missing verb being supplied from what preceded, i.e., “gives”. That this cannot be the meaning of \emph{dazu}, becomes clear from the Large Catechism and the Latin translations according to which \emph{dazu} is equivalent to “further”. This is the very circumstance that occasions the question as to the predicate of this part of the explanation whether it must be found in one of the preceding verbs or in one that follows, that is, in \emph{gegeben hat und noch erhaelt} or, perhaps, in \emph{versorgt}.

The last named construction enjoyed great vogue during the sixteenth century. We find it in the first Latin translation: “\emph{Credo, quod Deus me una cum omnibus aliis creaturis creaverit, quod mihi corpus et animam, hos oculos, aures, omnia alia membra, quod rationem et illos sensus donaverit. Neque id solum, sed credo quoque, quod omnia illa alias peritura sustentet, quod vestes, calceos, cibum et potum, certas sedes, uxorem, liberos, agros, iumenta et quidquid bonorum est, abundè ad sustentandam hanc vitam mihi impertiat praebeat. Credo etc.}” Again in the second translation: “\emph{Credo, quod Deus creavit me una cum omnibus creaturis: Quod corpus et animam, oculos, aures, et omnia membra, rationem et omnes sensus mihi dedit et adhuc sustentat. Ad haec quod vestes et calceos, cibum ac potum, domum, uxorem, liberos, agros,}
\quellseite{375}
\emph{iumenta cum omnibus vitae necessariis copiose et quotidie largitur: „Me contra omnia pericula protegit etc.”} Likewise in the third Latin translation: “\emph{Credo, quod Deus Pater me creavit . . . omnesque sensus dedit et conservat. Victum, vestitum, cibum, potum, domum, fundos, uxorem, liberos, iumenta et omnia bona, omnia necessaria ad vitam quotidie suppeditat, quod contra omnia pericula nos defendit, etc.}” Nor does the fourth Latin translation, by Hiob Magdeburg, show a different construction: “\emph{Credo, me cum tota rerum universitate a Deo creatum esse, qui corpus mihi et animam . . . sensusque omnes et largitus est et conservat. Vestibus mihi insuper et calciamentis, cibo et potu, aedibus, uxore et liberis, agris, pecoribus et ceteris bonis ad conservandam vitam necessariis abunde quotidie prospicit.}” Similarly in the Greek rendering in the Catechesis Graeco-Latina: “\emph{ἔτι τὰ ἐνδύματα καὶ τὰ φράγματα . . . καὶ πάντα τὰ ἄλλα ὑπάρχοντα . . . δίδωσιν}.” Accordingly, the words, \emph{dazu Kleider und Schuh . . . versorgt}, were accepted as a coherent statement, and the word \emph{versorgt} (=to provide someone with something) was taken in the sense of \emph{besorgt, schafft herbei} (=to procure something for someone), with the accusatives \emph{Kleider und Schuh} etc. as its object. \emph{Mit aller Leibsnahrung und Notdurft} would, accordingly, be detached from versorgt and become: “and, in addition, all that I need for this body and life”. The English translation corresponding most closely to this construction is found in Cranmer’s Catechism of 1548: “I believe that God the Father hath made me and all creatures in heaven and hearth, that he hath gyven to me and conserveth me bodye and soule, reason, senses, eyes, eares, and all my other members. And I believe that the same almightye Lord and God doth dayly gyve to me and to us all meat, dryncke, cloth, wife, children, house, lande, riches, cattell, and all thynges necessarye to the mayntenaunce of our lyves etc.” Joachim Moerlin (1547 or 1554) does not hesitate to paraphrase: “That God the Father daily and richly gives me these (namely body and life) and all things necessary to this life and temporal food”. The same construction is found again, e.g., in Marbach, 1557; in Aquila, 1538; in Bischoff, 1599; in Hercko, 1554; in Siber, 1575; in Dresser, 1581; in Hadamarius, 1552; in Maurer, 1575; in Joseph Opitz, 1538. The usage of versorgt as equivalent to besorgt was evidently not unknown to any of these; and because an example is found also in Luther (in a letter to Hans Schotten, May 1524, he says concerning the duties of the father: “\emph{er ist schuldig, dem Kind Essen und Trinken, Kleider . . . zu versorgen, für des Kindes Not und zu seinem Besten}”; cf. Erl. Ed. liii, p. 243),
\quellseite{376}
it might appear that this construction is the right one. It would, in that case, have the additional advantage of giving in the words: “I believe that God has made me my reason and all my senses” . . . a definition of the word “Creator,” and in the words “also clothing and shoes . . . guards and keeps me from all evil” an exposition of “Father almighty”. What an earthly father does for his children, Luther would raise to the highest degree and apply to the almighty Father.

Notwithstanding, this is hardly the correct construction. Taking the Small Catechism alone into consideration, there are three points which arouse our suspicion. (1) This construction would express that God richly and daily (!) provides me with wife and child; (2) it is incongruous, after “all goods”, an expression indicating that the enumeration is finished, to continue immediately with the words “and in addition, all that I need for this body and life”; and (3) even the rhythm, by no means unimportant in Luther’s explanations, requires a pause after “goods”. More significant is another factor: Luther’s preliminary labors militate against this construction. In the sermon of December 10, 1528, we read: “The meaning is that this I am to believe that I am God’s creature; that He has given me body, soul, sound eyes, reason, property, wife, children, fields, meadows, hogs, and cattle. Then that He has given the four elements etc. This article, accordingly, teaches that you do not have your life of yourself, not even a hair. All things in existence are comprehended in the little word Creator” (W. E. xxx 1, p. 87). In the Large Catechism we read: “What is the force of this, or what do you mean by these words: I believe in God, the Father almighty, Maker, etc.? Answer: This is what I mean and believe, that I am a creature of God; that is, that He has given and constantly preserves to me my body, soul, and life, members great and small, all my senses, reason, and understanding, and so on, food and drink, clothing and support, wife and children, domestics, house and home, etc. Besides, He causes all that exists to serve for the uses and necessities of life, sun, moon, and stars in the firmament, etc.” (W. E. xxx 1, p. 183; Tr. p. 681). These two quotations inevitably lead to the following conclusions: (1) That the section, “also clothing and shoes . . . all my goods” belongs to “he has given and still preserves to me” and that the objects here enumerated (though following the verb, in the German original) are to be taken as coordinate with the first named objects (body, soul, etc.). (2) That there is by no means to be supplied, as, for instance, Luehrs and Zezschwitz require, a \emph{mich} in connection with \emph{erhaelt} which would produce the
\quellseite{377}
following text: “that He has given me my body and soul, eyes, ears, and all my members, my reason, and all my senses, and still preserves me”; but that the \emph{mir} belongs both to \emph{gegeben hat} and to \emph{erhaelt}, thus: “that He has given and still preserves to me my body and soul etc.” This, then, is the thought expressed: I believe that God has given and still preserves to me body and soul, eyes, ears . . . . all senses, further, clothing, shoes, meat, drink, . . . and all goods; and that He also preserves to me all these objects. In preserving them He, indeed, preserves also me; but by interpolating mich as the object of erhalt we would destroy the whole sentence structure that Luther intended to rear. (3) That the words mit aller Leibesnahrung und Notdurft . . . versorgt constitute a new, complete thought, in the wording of which the terms \emph{versorgt} and \emph{mit} are to be taken in the usual sense, instead of being exchanged for besorget and nebst. The absence of a direct mich does not disprove this statement, for Luther often omits the personal pronoun in the second part of a sentence, and does this even if the cases do not agree (cf. the sentence in the Large Catechism which expresses the very same thought: “\emph{Da sehen wir, wie sich der Vater uns gegeben hat samt allen Kreaturen und aufs allerreichlichste in diesem Leben versorgt}, Tr. p. 682, § 24). (4) That the conventional interpretation of \emph{Notdurft und Nahrung dieses Leibes und Lebens} according to which these words are a summary of the blessings previously enumerated, is an error (even Kaftan still adheres to this error although in other respects he has consistently rejected the construction of the Latin translations). Luther, far from virtually repeating his statement, thinks of every necessity of this life not previously named, as air and light, fire and water, sun, moon, stars, etc. Careful reading of the Large Catechism shows that Luther in this phrase includes everything enumerated by him in the explanation of the Fourth Petition alongside of “clothing, shoes . . . all goods.” That thus, especially through his reference to the elements, a vital and timely thought is introduced, is evident to everyone who knows how often in our day nature is placed alongside of God, if not above Him, as something independent (cf. p. 315). Compare the lines of Paul Gerhardt: “Heaven, earth, and all their host, As my servants He has given, etc.” When we remember that Luther had previously made mention of the “creatures”, and when we have come to realize that the words “has given my body and sould” down to “richly and daily provides me with all that I need for this body and life,” is nothing but an unfolding of the preceding main clause: “I believe that God has made me and all
\quellseite{378}
creatures”, there will be still less desire to eliminate his reference to air, light, etc. (5) From the Catechetical Sermon quoted above (p. 354) and the Large Catechism it is clear that the giving, preserving, providing, here mentioned are by Luther intended to serve as a definition of divine creation. For this reason it is totally wrong to add to creation “preservation” which is said to express itself in giving (clothing, shoes, etc.), providing and protecting (so Opitz, H. Magdeburg); and even greater is the error when also “government” is intercalated. The benefit (?) resulting from this procedure is that we have in happy union the dogmatical trio—creation, preservation, government. “All is comprehended in the word creator. Emphasize the word creator!” says Luther. In the Small Catechism he even goes beyond his preliminary labors in that he so extends the statement: “I believe that God has made me etc.”, as to find in it the notion of protecting, guarding, and keeping. This has been recognized, e.g., in the Joachimstal Catechism (1574) where we read: “Creation does not merely mean that God once upon a time made everything but that He is present with His creature and active for its sake, as the words of the explanation (down to “from all evil”) and the hymn (Luther’s hymn on the Creed) show” (Reu i 2, p. 696). All this—giving, providing, protecting, etc., can, according to the Large Catechism, be classified under a double head: (1) God has given us everything we have and see before us; (2) He also wards off whatever could hurt me and mine.

The Catechetical Sermons of 1528 and the Large Catechism justify the inference that all the statements in the Small Catechism are designed to lay stress upon God’s activity as contrasted with human wisdom and strength. In the Sermon of December 10 we read: “This article teaches, accordingly, that you do not have life of yourself, not even a hair . . . Whatever you may possess, no matter how little it be, remember this when you say ‘Maker.’ Even though you wear a wreath on your head—let us not think that we have made ourselves, as do the proud princes . . . Nothing of all this I have of myself. The Maker, that is, God, has given everything . . . If all things are God’s gifts, you owe it to Him to praise and thank Him for the same . . . How many, I ask, are in the world who understand this word ‘Creator’? For no one serves Him. Accordingly the First Article may well humble and dismay us because we do not believe . . . He, however, who believes this article, when he looks at his cow, will say, ‘The Lord has given it to me.’ The same in regard to his wife, his children . . . I have nothing of
\quellseite{379}
myself lest I should become proud. I can neither give it to me nor preserve it” (W. E. xxx 1, p. 87 f.). Compare also the Nuremberg Sermons for Children (Reu i 1, p. 510). This, then, is the construction of the first part of the explanation: I believe that it is God who made me together with all creatures; that is (1) I believe that whatever I am and have is from Him alone; (a) I believe that it is God who has given me my body and soul; (b) I believe that it is God alone who also preserves to me my body and soul; (c) I believe that whatever else I have besides body and soul is given me and preserved for me by Him alone; (d) I believe that it is God alone who also provides me with all I need for this body and life (for, being Lord of the world, He gives me rain and sunshine, light and air, fire and water; yes, He has appointed heaven and earth, air and sea for my service); (2) I believe that it is God who wards off whatever could hurt me and mine; (a) I believe that it is He alone who protects me and mine from all danger, (b) guards me from all evil, (c) keeps me from all evil. Should the argument be urged against this construction that the “me” it “provides me with all that I need for this body and life” is an intercalation (\emph{mit aller Notdurf und Nahrung dieses Leibes und Lebens reichlich und taeglich versorgt}), let it be remembered that the other construction has likewise an intercalation, in that the “me” in “keeps me from all evil” is not in the original (\emph{wider alle Faehrlichkeit beschirmet}); and that Luther has by no means always been exact in the use of the personal pronoun so far as the Small Catechism is concerned. A case in point is the second half of the explanation of the Fifth Commandment or the third part of the explanation of the First Article. It is, accordingly, quite likely that Luther when he used the verbs “protect, guard, keep”, had in mind as object not only “me” but every object enumerated in the preceding explanation. For this reason we have explained in the epitome above given: “Me and mine” (see Large Catechism, Jacobs, Book of Concord, p. 440, § 17; Tr. p. 680 adds uns). On the whole question of construction compare Reu, \emph{Die Konstruktion des ersten Artikels} in \emph{Kirchl. Ztschr.}, 1904, pp. 28–33; 83–89.

Finally we must note how well Luther succeeds in filling the article of Creation with contents of vital significance for the believer. He attained this end principally by relating God’s creative activity to the individual (cf. p. 105), but also by viewing all creatures as existing for our benefit (e.g. in “provides me with all that I need
\quellseite{380}
etc.”); and the fullness of assurance with which Luther confesses that God protects, guards, and keeps us from all evil, has its basis in just this faith that God is the maker of all creatures: He is the Lord of all, no creature can thwart His plan, all of them are bound to be His instruments in protecting, guarding, keeping His children. Magnificent is Luther’s presentation of the subject elsewhere (Walch viii, p. 318): “If we should behold the downfall and collapse of the world and all its elements, and ourselves floundering in the wreck, we still should say: You will not fall though you do fall unless it be the will of God. And though the world should lie, a huge weight, upon our very head, we still should say: World, you can not hurt or crush me; but if it should please God to have me felled and crushed by your weight, let there be done in the name of the Lord whatever pleases Him: my time is in His hands. But should the contrary be His design, I will defy you, heaven and earth, and all power of pope and Turk and the whole world’s wrath with it” (cf. p. 296). The same thought he expresses in the Short Form: “If He is almighty, what can I lack that He cannot give me and do for me? If He is Creator of heaven and earth and Lord of all things, who will take anything from me, or harm me?” (above p. 85). Not all of the catechists of the sixteenth century have been unmindful of the fact here brought out. Compare the Nuremberg Sermons for Children (Reu i 1, p. 405), and what Meder (1590) says on basis of these sermons: “If God has made heaven and earth, He is also Lord of them; and everything in heaven and earth must move according to His will; and if He has made everything for His sake, everything is bound to serve us” (Reu i 1, p. 495).

For lack of space, the remainder has to be compressed into a few brief notes; for a more detailed discussion see M. Reu, \emph{A New English Translation of Luther’s Small Catechism}, 1926. In the text of the Second Article, the edition of 1531 emphasizes the word “Lord” as the one most important by capitalizing every letter: HERR. In keeping with this conception, clearly expressed in the sermon of December 10, 1528, and in the Large Catechism where that sermon is largely reproduced, Luther chooses the heading “Of Redemption”, and groups the whole explanation around the statement, “I believe that Jesus is my Lord”. Already at an early day, this simple and lucid construction has unfortunately been misunderstood by some (for instance, H. Magdeburg: “\emph{Credo Jesum Christum verum}
\quellseite{381}
\emph{esse Deum et a patre ab æterno esse genitum: veritum item ex Maria virgine hominem natum, eundemque dominum meum.}” Still worse is W. Han’s product, of 1553; Reu i 2, 2, p. 583; i. 2, 1, p. 459); later it was ignored almost universally in the interest of dogmatics. If there is any point in regard to which the Large Catechism, here a virtual reproduction of that sermon of December 10, 1528, can serve as commentary upon the Small Catechism, it is the one under consideration; and for this reason everyone should study the Large Catechism who would grasp Luther’s intentions and escape the danger of mixing foreign ingredients with his Catechism.—“Redeemed, purchased and won” are words expressing one and the same thought; but each time under a different image. The first of these words reminds one of the great price paid for us; the second of the arduous labor necessary; the third of the fierce battle that was fought.—The section stating the ultimate purpose of redemption (“that I might be His own, etc.”) treats solely of God’s, of Christ’s, purpose, and should not be deprived of its objective character. It remains for the Third Article to picture the execution of this purpose in detail; there the very question is answered, how the individual is brought to the enjoyment of all the blessings which Christ has won for him by His redemptive work. Instead of anticipating the Third Article, the object must be to let the pupil envisage the blessedness of the change purposed by Christ (formerly the devil’s property, now Christ’s; formerly a lost and condemned sinner, now righteous, innocent, and blessed). This blessedness is to begin here below (“in His kingdom” describes the life of the faithful in the kingdom of Christ as it is here on earth, i.e., the Church; cf. W. E. xxx 1, pp. 80 f.; 186); it includes, however, the life beyond (cf. “in everlasting righteousness”, a phrase equivalent to “state of justification”, cf. W. E. xxx 1, p. 187). It is to be hoped that no catechist will henceforth dissect this magnificent period and endeavor to correlate its constituent parts with this or that part of the text of the Apostolic Creed. He would deprive “the golden gem” of the fairest part of its beauty and dim its power (cf. pp. 103–105).

He who has come to the Third Article by way of the First and Second, will expect Luther to select some one word as the nucleus for the whole explanation. The heading, “Of Sanctification”, strengthens this expectation. And Luther actually begins this article in his Large Catechism: “This article I cannot relate better than to Sanctification, that through the same the Holy Ghost, with His
\quellseite{382}
office is declared and depicted, namely, that He makes holy. Therefore we must take our stand upon the word Holy Ghost, because it is so precise and comprehensive that we cannot find another” (Tr. p. 687 § 35). But how does Luther now correlate the other parts of this article with the main one? The Large Catechism does not leave us without an answer to this question. We read presently: “But how is such sanctifying done? Answer. Just as the Son obtains dominion, whereby He wins us, through His birth, death, resurrection, etc., so also the Holy Ghost effects our sanctification by the following parts, namely by the communion of saints or the Christian Church, the forgiveness of sins, the resurrection of the body, and the life everlasting; that is, He first leads us into His holy congregation, and places us in the bosom of the Church, whereby He preaches to us and brings us to Christ” (Tr. p. 687 § 37). Subsequently he repeats this outline of the article in practically identical words and says of the exposition given on basis of this outline: “Let this suffice concerning the sum of this article.” But he himself does not let this “sum” suffice. “Rather, because the parts which are here enumerated are not quite clear to the simple, we shall run over them also”, he says (Tr. p. 689 § 46), and then devotes a special paragraph each to Church, Communion of saints, Forgiveness of sins, Resurrection of the body, and Life everlasting. Not so in the Small Catechism. Here he consistently carries out his principle that “everything must take its stand upon the word Holy Ghost.” He touches, indeed, upon the other parts, but he neither coordinates them with the main part nor explains them as \emph{disiecta membra}; he subordinates them all to the Holy Spirit, and confesses his faith in Him and His work alone. At a previous place we have emphasized the fact that these attributes of unity and compactness constitute the principal excellence of his explanation of the Third Article (p. 105).

It is regrettable that this advantage has been recognized and turned to account by but few commentators. Quite a number have adopted the mediæval twelvefold division, while others, who succeeded in preserving the unity of the first two articles, have dissected the Third into five, more or less discrete parts. This has been the case down to modern times when the assertion was made, or rather renewed, that Luther had accurately followed in his explanation the several parts of the text. There is some truth in this; but when the subjects, “Church”, “Forgiveness” etc., are treated independently, the fact is overlooked
\quellseite{383}
that Luther does not see them at all in that light, but only as subordinate to the one subject—the Holy Spirit. Even Kaftan, although more intent than most catechists upon maintaining the connection between the Holy Spirit as subject and the other parts, nor unmindful of the tense employed by Luther (“has called me . . . kept me in the true faith”), does not do justice to the matter. He says in regard to the division: “Besides our faith in the Holy Ghost we also confess in the Third Article that there is a holy Christian Church, which is His workmanship and in which He Himself dwells and works” (Kaftan-Horine, p. 131). Now let us apply to this statement the touch-stone of Luther’s exposition: according to that, we confess the Church not as something “alongside the Holy Spirit”, but we confess of the same Holy Spirit, in regard to whom we have confessed something in the first section, something additional, namely, that He calls the whole Christian Church on earth, etc. We have been of the conviction for more than thirty years that Luther’s explanation best comes into its own when outlined as follows: (1) What the Holy Spirit has done in the past; (2) what He is doing in the present; (3) what He still will do in the future. When the article is thus outlined, everything, precisely as in Luther’s explanation, is \emph{dominated by one subject}—the Holy Spirit; the other parts of the article receive the same treatment that Luther accorded them; the perfect tense in the first section which is so largely but none the less erroneously exchanged for a present tense, retains its distinct tense value; and the whole article, as we can assert from experience, becomes clear and intelligible to the children. Nor is it necessary to repeat here that, where the Third Article is thus taught, the fundamental fact that the salvation of the Christian, from beginning to end, rests not on himself but on Christ and His Holy Spirit, will effectively be impressed upon the pupil’s consciousness (pp. 105 f.). Kaftan, briefly, repudiates this outline (p. 229). But the Large Catechism shows that it was not at all remote from Luther’s mind; for he there connects the subjects of Church and forgiveness of sins with the present, and those of the resurrection and eternal life with the future, exactly as he does in the Small Catechism. Here are his words: “Therefore we believe in Him who through the Word daily brings us into the fellowship of this Christian Church, and through the same Word and the forgiveness of sins bestows, increases, and strengthens faith, in order that when He has accomplished it all, and we abide therein, and die to the world and to all evil, He may finally make us perfectly and forever holy; which now we expect
\quellseite{384}
in faith through the Word” (W. E. xxx 1, p. 191; Tr. p. 695 § 62). With even greater clearness, these thoughts are found in the Catechetical Sermon of December 10, 1528: “He is here called the Holy Spirit. Why? Because He sanctifies. And I believe in the Holy Spirit for the reason that He has sanctified me and sanctifies me to this day . . . And the Holy Spirit will sanctify you through the resurrection of the body . . . then the Holy Spirit will complete His work and my sanctification will be complete . . . The Third Article means that I believe in the Holy Spirit, that is to say, that the Holy Spirit of God will and does sanctify me . . . He has already begun to sanctify me; when I shall have died, He will complete this sanctification through these two parts, the resurrection of the body and eternal life . . . The other parts signify the means, or method, through which He sanctifies me, for the Holy Spirit does not sanctify you outside the Church. That is the reason that the Christian Church in which all His gifts are found, is placed immediately after the Holy Spirit. Through this Church He proclaims His Message, calls you, makes Christ known to you, and kindles faith in you that, through the Word and the Sacraments, you may become free from sins; and thus you are truly free on earth. If you die and remain in the Church, He will raise you from the dead and will sanctify you wholly . . . And the Holy Spirit sanctifies me through the Word and Sacraments, which are in the Church, and will sanctify us perfectly on the last day” (W. E. xxx 1, pp. 91. 93 f.). We refer to our Explanation as proof that when the Third Article is thus skeletonized, Luther’s explanation comes into its own better than with any other outline so far suggested.

The term sanctification occurs in the Third Article in both a narrower and wider sense: in the wider sense in the heading, in the narrower in the sections “sanctified in the true faith” and “sanctifies the whole Christian Church”. Sanctify in the wider sense means to Luther “to lead to Christ for time and eternity”; compare the above quotations. In the narrower sense, therefore, no other meaning is possible than “to bring one to faith and thus to lead him into a state of righteousness.” The latter view has been impugned by those who understand by “sanctified” subjective sanctification—the victorious life-long conflict with sin, or at least, the union of both—justification and the daily subjective sanctification. That sanctification can be taken in a subjective sense, and has been taken in that sense also by Luther, is perfectly true; but that
\quellseite{385}
does not prove that he used it in that sense in this connection. That Luther has rather used it at this place in an objective sense, is proved in the first place by the fact that the objective sense, in his writings as well as in the New Testament, is the one which predominates (cf. Koestlin, \emph{Luther’s Theologie} ii, pp. 178 f. “sanctification = justification”); in the second place, by the fact that it is so used in the Large Catechism, e.g., “Therefore sanctifying is nothing else than bringing us to Christ to receive this good [Christ’s redemption] to which we could not attain of ourselves” (W. E. xxx 1, p. 188; Tr. p. 689 § 39).—With this question another one is intimately connected, namely, how the four verbs “\\emph{call, enlighten, sanctify, and preserve}” are to be understood. The view has found favor that they denote four stages of the \emph{ordo salutis} and that they must be interpreted in strictest conformity to our dogmatical terminology (some have found even a chronological order in these words!). Thus e.g., Tetelbach: “How many acts of the Holy Spirit are here enumerated? Four. First, He calls us to a diligent hearing of the Word of God, and makes and effects in us the beginning of conversion. Which is the other? He enlightens by the gospel our dark and blind hearts through the light of faith so that we, erstwhile unbelievers, become believers. Which is the third? He sanctifies us in the true faith so that a new life proceeds from faith, a life of new obedience, of love, joy, peace, patience, kindness, charity, meekness, and chastity. Which is the fourth? He preserves us constantly in the truth which we have recognized, in the true faith, and a blessed confession to the end. Accordingly, beginning, middle, and end of our salvation is the work of God the Holy Spirit” (Reu i 1, p. 689 f.). Formerly we had the same view; but by paying more attention to biblical usage regarding this point (see Reu, \emph{Heilsordnung}, 1914), and by renewed study of Luther’s preliminary labors, we have come to the conclusion that the first three of the verbs enumerated are not at all intended to describe three distinct, successive stages of the Holy Spirit’s work, but rather the one saving work from three different aspects, in order effectively to reveal its magnitude. When we were far from Christ, the Holy Spirit led us to Christ (call—effectual, successful call, not invitation); when darkness was in our hearts, He enlightened us with His gifts, i.e., He kindled a bright light in our hearts by the knowledge which He imparted; when we were servants of sin, He transplanted us into a state of holiness and righteousness, through faith in the forgiveness obtained in Christ, a faith wrought
\quellseite{386}
by Himself. In the first clause the emphasis is on the Word through which the act of sanctification takes place; in the second, on the knowledge conveyed by the Holy Spirit; in the third, on saving faith as wrought by the Holy Spirit. Looking back, the Christian cannot but confess: “If the Holy Spirit had not led me to Christ through the gospel, I should still be far from Him; if He had not enlightened my heart with His gift of knowledge, I should still sit in darkness; if He had not kindled faith in Christ in my heart and thereby enabled me to obtain forgiveness of sins and to enter into a state of holiness and righteousness, I should still be in my sin”. But in making this confession, the Christian does not enumerate three successive acts; he merely means to give expression, from every possible point of view, to the one blessed experience of salvation. This was understood by A. Siber already who paraphrased this section of the Enchiridion in his \emph{Sabbatum puerile} (1575) as follows: “\emph{Qui quidem, lege conclusum et occisum, sermone reconciliationis, quae est vox Evangelii, in vitam me vocavit} (note well: effectual call!), \emph{sua gratia collustravit, fidei sanctitate insigniverit}” (Reu i 2, 2, p. 570). When discussing this section with its past tense, the catechist will have to impress upon the children the fact that the Holy Spirit has worked upon them through Baptism and instruction in the Word; but he dare not forget what was said by us on p. 314: let him speak from the standpoint of his own life or that of well-known men of God lest guile or lip-service insinuate themselves into the teaching, especially in view of the fact that he also confesses: “and preserved me in the true faith.”

Finally, a brief remark relative to the much mooted section, “enlightened me with His gifts”. He who thinks that the fundamental features, or the various successive steps, of the order of salvation are here indicated, and accordingly mistakes the “call” as a mere invitation to come to Christ, will surely find himself tempted to understand by the “gifts” with which he has been enlightened, the Law and the Gospel—a view that, for a long time, constituted the exclusive catechetical tradition. While already the sequence thus arising (Gospel, Law, Gospel) should make one suspicious, that it in itself would not be an invincible obstacle, at least not from the standpoint of dogmatics. Even more significant would be the change of prepositions: for we should expect durch, by, a second time instead of mit, with. What is more important is the fact that nothing in Luther justifies such an interpretation, at least nothing in his preliminary labors on the Catechism. Albrecht
\quellseite{387}
has well summed up what is to be said in this respect: “In the sermon of December 10 (W. E. xxx 1, p. 94), the \emph{dona spiritus sancti} evidently mean the objective spiritual gifts or official endowments in the sense of 1 Cor. 12:4 ff. and Eph. 4:11 f. The phrase “all His gifts” in the Large Catechism will have to be understood in the same sense (W. E. xxx 1, p. 192), probably also that found on page 190: “with various gifts.” But the clause on page 192, line 15 f.: “and besides they are not enlightened and favored by any gifts of the Holy Spirit”, when studied in the light of its context, leads to the conclusion that here the subjective operations of the Holy Spirit are meant. Luther’s point of view may have been the widely spread medieval doctrine of the seven gifts of the Holy Spirit, according to Isa. 11:2 (that Luther retained this view is proved from his pentecostal hymn, “Come, O Creator God . . . . Thou art with powers sevenfold the finger on the hand divine); or it may have been certain biblical passages, such as Rom. 8:1 ff.; Gal. 5:22; or possibly a suggestion such as that contained in the German Patrem: \emph{Der aller Bloeden Troester heisst und mit Gaben zieret schoene}” (W. E. xxx 1, p. 367 f.). R. Hadamarius writes (1552): “\emph{Quae dona sunt Spiritus sancti? Vera sapientia, intellectus, consilium, fortitudo, scientia, pietas, timor Dei}” (Reu i 2, 2, p. 1030). Bischoff (1599): “Thereupon He enlightens our hearts with the true knowledge of God and true faith in Jesus Christ, whereby we obtain the forgiveness of sin” (Reu i 2, 2, p. 231). The Joachimstal Catechism (1574): “Who enlightens our blind hearts? Alone the Holy Spirit, the very love and fire in God; He is poured out into our hearts and regenerates us, enkindling a new light in us and leading us to the knowledge of God the Father and His Son, Matt. 16:17; 1 Cor. 12:3” (Reu i 2, 2, p. 609). The interpretation of the gifts of the Spirit in an objective sense did not appear in the expositions of the sixteenth century. In itself it would be quite correct to say that the Holy Spirit enlightens us by arousing apostles, prophets, evangelists, etc. A warrant for this view might be found in the December sermons and the stress laid upon the ministerial office in the Large Catechism; but this interpretation, too, would call for the preposition durch, by, instead of mit, with. The subjective gifts must be meant. The Spirit has enlightened and favored, graced and adorned me with knowledge, wisdom, comfort, power, and every variety of virtue (cf. Large Catechism, W. E. xxx 1, pp. 190, 192; Tr. p. 691 § 51, 697 § 66); or, as more in harmony with the context in the Small Catechism (consider, e.g., the contrast to
\quellseite{388}
“not by my own reason”): He has enlightened me within by giving me the true knowledge of Christ and my sin: of my sin, by teaching me that it poisons my life and accuses me before God; of Christ, by showing me that He is my Saviour and my Lord, who has redeemed me from sin, and whose own I am to be henceforth and forever.

Whatever needs to be said concerning the \emph{Holy Trinity}, will readily suggest itself when the teacher at the conclusion of his discussion sums up the contents of the three Articles. “We could never attain to the knowledge of the grace and favor of the Father except through the Lord Christ, who is a mirror of the paternal heart, outside of whom we see nothing but an angry and terrible Judge. But of Christ we could know nothing either, unless it had been revealed by the Holy Ghost”, says Luther in the Large Catechism (Tr. p. 695 § 65).—In regard to the other Chief Parts we refer to Kaftan and our \emph{Explanation}, cf. also M. Reu, \emph{A New Translation} etc.

Whether it be advisable to supply the child with an elaborate explanation of the Small Catechism as an aid in instruction is a moot point. Our local educational conditions would strongly speak in its favor. In any event, the printed explanation is to be considered only as \emph{the result or thesis of the previous catechesis}; it is given to the child in order that he may have a fixed lesson for home study. For the catechesis proper pupil as well as teacher knows only the text of Luther’s Catechism. Whether the form of such catechetical explanation is erotematic or thetic, is in itself immaterial. This much, however, can be said in favor of the thetic form: (1) the catechist will find himself less \emph{shackled} by the employment of the thetic form than the other; (2) the child is more easily trained for independent thinking where the thetic method obtains for the reason that, in reviewing, he is compelled to seek the answer for himself; (3) the thetic form supplies the thought as a whole whereas the other method usually supplies only half the thought, since the child is very much inclined to memorize only the answer and to put the question aside; and (4) the thetic method leads to better mastery of the Catechism, for here the text of the Catechism can be plainly traced by the eye, as it winds through all the explanation like a string of pearls.
